%%%%%%%%%%%%%%%%%%%%%%%%%%%%%%%%%%%%%%%%%%%%%%%%%%%
%% LaTeX book template                           %%
%% Author:  Amber Jain (http://amberj.devio.us/) %%
%% License: ISC license                          %%
%%%%%%%%%%%%%%%%%%%%%%%%%%%%%%%%%%%%%%%%%%%%%%%%%%%

\documentclass[a4paper,11pt]{book}
\usepackage[T1]{fontenc}
\usepackage[utf8]{inputenc}
\usepackage{lmodern}
%%%%%%%%%%%%%%%%%%%%%%%%%%%%%%%%%%%%%%%%%%%%%%%%%%%%%%%%%
% Source: http://en.wikibooks.org/wiki/LaTeX/Hyperlinks %
%%%%%%%%%%%%%%%%%%%%%%%%%%%%%%%%%%%%%%%%%%%%%%%%%%%%%%%%%
\usepackage{hyperref}
\usepackage{graphicx}
\graphicspath{ {imagens/} } 
\usepackage[portuguese]{babel}
\usepackage{booktabs}
\usepackage[table,xcdraw]{xcolor}
\usepackage{textcomp}
\usepackage{booktabs}
\usepackage{tikz}
\usepackage{pbox}
\hypersetup{
	colorlinks,
	linkcolor={red!50!black},
	citecolor={blue!50!black},
	urlcolor={blue!80!black}
}

\usepackage{eso-pic}
\newcommand\BackgroundPic{%
	\put(0,0){%
		\parbox[b][\paperheight]{\paperwidth}{%
			\vfill
			\centering
			\includegraphics[height=\paperheight]{fundo}
			%\includegraphics[width=0.5\paperwidth,height=\paperheight,keepaspectratio]{imagens/processador}%
			%\includegraphics[width=\paperwidth,height=\paperheight]{imagens/fundo.png}%

			%\vfill
}
{\put(-320,-10){\makebox[0cm][c]{\raisebox{9cm}{\includegraphics[height=4cm]{processador}}
}}}}}

%%%%%%%%%%%%%%%%%%%%%%%%%%%%%%%%%%%%%%%%%%%%%%%%%%%%%%%%%%%%%%%%%%%%%%%%%%%%%%%%
% 'dedication' environment: To add a dedication paragraph at the start of book %
% Source: http://www.tug.org/pipermail/texhax/2010-June/015184.html            %
%%%%%%%%%%%%%%%%%%%%%%%%%%%%%%%%%%%%%%%%%%%%%%%%%%%%%%%%%%%%%%%%%%%%%%%%%%%%%%%%
\newenvironment{dedication}
{
   \cleardoublepage
   \thispagestyle{empty}
   \vspace*{\stretch{1}}
   \hfill\begin{minipage}[t]{0.66\textwidth}
   \raggedright
}
{
   \end{minipage}
   \vspace*{\stretch{3}}
   \clearpage
}

%%%%%%%%%%%%%%%%%%%%%%%%%%%%%%%%%%%%%%%%%%%%%%%%
% Chapter quote at the start of chapter        %
% Source: http://tex.stackexchange.com/a/53380 %
%%%%%%%%%%%%%%%%%%%%%%%%%%%%%%%%%%%%%%%%%%%%%%%%
\makeatletter
\renewcommand{\@chapapp}{}% Not necessary...
\newenvironment{chapquote}[2][3em]
  {\setlength{\@tempdima}{#1}%
   \def\chapquote@author{#2}%
   \parshape 1 \@tempdima \dimexpr\textwidth-4\@tempdima\relax%
   \itshape}
  {\par\normalfont\hfill--\ \chapquote@author\hspace*{\@tempdima}\par\bigskip}
\makeatother

%%%%%%%%%%%%%%%%%%%%%%%%%%%%%%%%%%%%%%%%%%%%%%%%%%%
% First page of book which contains 'stuff' like: %
%  - Book title, subtitle                         %
%  - Book author name                             %
%%%%%%%%%%%%%%%%%%%%%%%%%%%%%%%%%%%%%%%%%%%%%%%%%%%
% Book's title and subtitle
\title{\Huge \textbf{x86-64 \\ \Huge Programação Assembly com Ubuntu}}
% Author
\author{\textsc{Ed Jorgensen}\\ {Tradução e adaptação: João Araujo}\thanks{araujo at eng.uerj.br}}

\usetikzlibrary{shapes,arrows}

% taken from manual
\makeatletter
\pgfdeclareshape{document}{
	\inheritsavedanchors[from=rectangle] % this is nearly a rectangle
	\inheritanchorborder[from=rectangle]
	\inheritanchor[from=rectangle]{center}
	\inheritanchor[from=rectangle]{north}
	\inheritanchor[from=rectangle]{south}
	\inheritanchor[from=rectangle]{west}
	\inheritanchor[from=rectangle]{east}
	% ... and possibly more
	\backgroundpath{% this is new
		% store lower right in xa/ya and upper right in xb/yb
		\southwest \pgf@xa=\pgf@x \pgf@ya=\pgf@y
		\northeast \pgf@xb=\pgf@x \pgf@yb=\pgf@y
		% compute corner of ‘‘flipped page’’
		\pgf@xc=\pgf@xb \advance\pgf@xc by-10pt % this should be a parameter
		\pgf@yc=\pgf@yb \advance\pgf@yc by-10pt
		% construct main path
		\pgfpathmoveto{\pgfpoint{\pgf@xa}{\pgf@ya}}
		\pgfpathlineto{\pgfpoint{\pgf@xa}{\pgf@yb}}
		\pgfpathlineto{\pgfpoint{\pgf@xc}{\pgf@yb}}
		\pgfpathlineto{\pgfpoint{\pgf@xb}{\pgf@yc}}
		\pgfpathlineto{\pgfpoint{\pgf@xb}{\pgf@ya}}
		\pgfpathclose
		% add little corner
		\pgfpathmoveto{\pgfpoint{\pgf@xc}{\pgf@yb}}
		\pgfpathlineto{\pgfpoint{\pgf@xc}{\pgf@yc}}
		\pgfpathlineto{\pgfpoint{\pgf@xb}{\pgf@yc}}
		\pgfpathlineto{\pgfpoint{\pgf@xc}{\pgf@yc}}
	}
}
\makeatother


\begin{document}
\AddToShipoutPicture*{\BackgroundPic}

\frontmatter
\maketitle

\parskip=10.0mm
\noindent
\textbf{Cover image:}\\
Top view of an Intel central processing unit Core i7 Skylake type core,	model 6700K, released in June 2015.\\
Source: Eric Gaba, \url{https://commons.wikimedia.org/wiki/File:Intel\_CPU\_Core\_i7\_6700K\_Skylake\_top.jpg}

\noindent
\textbf{Cover background:}\\
By Benjamint444 (Own work)\\
Source: \url{http://commons.wikimedia.org/wiki/File%3ASwirly_belt444.jpg} 

\noindent	
Copyright 2015, 2016, 2017, 2018, 2019 por Ed Jorgensen\\
Copyright da tradução 2020 por João Araujo

\includegraphics[width=5cm]{imagens/assembly-007}


\parskip=2.0mm
\noindent
\textbf{Você é livre para:}\\
\textbf{Compartilhar} -- copiar e redistribuir o material em qualquer suporte ou formato.\\
\textbf{Adaptar} -- remixar, transformar, e criar a partir do material

\noindent
\textbf{De acordo com os termos seguintes:}\\
\textbf{Atribuição} -- Você deve dar o crédito apropriado, prover um link para a licença e indicar se mudanças foram feitas. Você deve fazê-lo em qualquer circunstância razoável, mas de nenhuma maneira que sugira que o licenciante apoia você ou o seu uso.\\
\textbf{Não Comercial} -- Você não pode usar o material para fins comerciais.\\
\textbf{Compartilha Igual} -- Se você remixar, transformar, ou criar a partir do material, tem de distribuir as suas contribuições sob a mesma licença que o original.

%%%%%%%%%%%%%%%%%%%%%%%%%%%%%%%%%%%%%%%%%%%%%%%%%%%%%%%%%%%%%%%
% Add a dedication paragraph to dedicate your book to someone %
%%%%%%%%%%%%%%%%%%%%%%%%%%%%%%%%%%%%%%%%%%%%%%%%%%%%%%%%%%%%%%%
%\begin{dedication}
%Dedicated to Calvin and Hobbes.
%\end{dedication}

%%%%%%%%%%%%%%%%%%%%%%%%%%%%%%%%%%%%%%%%%%%%%%%%%%%%%%%%%%%%%%%%%%%%%%%%
% Auto-generated table of contents, list of figures and list of tables %
%%%%%%%%%%%%%%%%%%%%%%%%%%%%%%%%%%%%%%%%%%%%%%%%%%%%%%%%%%%%%%%%%%%%%%%%
\tableofcontents
\listoffigures
\listoftables

\mainmatter

%%%%%%%%%%%
% Preface %
%%%%%%%%%%%
\chapter*{Preface}
Lorem ipsum dolor sit amet, consectetur adipiscing elit. Duis risus ante, auctor et pulvinar non, posuere ac lacus. Praesent egestas nisi id metus rhoncus ac lobortis sem hendrerit. Etiam et sapien eget lectus interdum posuere sit amet ac urna.

\section*{Un-numbered sample section}
Lorem ipsum dolor sit amet, consectetur adipiscing elit. Duis risus ante, auctor et pulvinar non, posuere ac lacus. Praesent egestas nisi id metus rhoncus ac lobortis sem hendrerit. Etiam et sapien eget lectus interdum posuere sit amet ac urna. Aliquam pellentesque imperdiet erat, eget consectetur felis malesuada quis. Pellentesque sollicitudin, odio sed dapibus eleifend, magna sem luctus turpis.

\section*{Another sample section}
Lorem ipsum dolor sit amet, consectetur adipiscing elit. Duis risus ante, auctor et pulvinar non, posuere ac lacus. Praesent egestas nisi id metus rhoncus ac lobortis sem hendrerit. Etiam et sapien eget lectus interdum posuere sit amet ac urna. Aliquam pellentesque imperdiet erat, eget consectetur felis malesuada quis. Pellentesque sollicitudin, odio sed dapibus eleifend, magna sem luctus turpis, id aliquam felis dolor eu diam. Etiam ullamcorper, nunc a accumsan adipiscing, turpis odio bibendum erat, id convallis magna eros nec metus.

\section*{Structure of book}
% You might want to add short description about each chapter in this book.
Each unit will focus on <SOMETHING>.

\section*{About the companion website}
The website\footnote{\url{https://github.com/amberj/latex-book-template}} for this file contains:
\begin{itemize}
  \item A link to (freely downlodable) latest version of this document.
  \item Link to download LaTeX source for this document.
  \item Miscellaneous material (e.g. suggested readings etc).
\end{itemize}

%%%%%%%%%%%%%%%%%%%%%%%%%%%%%%%%%%%%
% Give credit where credit is due. %
% Say thanks!                      %
%%%%%%%%%%%%%%%%%%%%%%%%%%%%%%%%%%%%
\section*{Acknowledgements}
\begin{itemize}
\item A special word of thanks goes to Professor Don Knuth\footnote{\url{http://www-cs-faculty.stanford.edu/~uno/}} (for \TeX{}) and Leslie Lamport\footnote{\url{http://www.lamport.org/}} (for \LaTeX{}).
\item I'll also like to thank Gummi\footnote{\url{http://gummi.midnightcoding.org/}} developers and LaTeXila\footnote{\url{http://projects.gnome.org/latexila/}} development team for their awesome \LaTeX{} editors.
\item I'm deeply indebted my parents, colleagues and friends for their support and encouragement.
\end{itemize}
\mbox{}\\
%\mbox{}\\
\noindent Amber Jain \\
\noindent \url{http://amberj.devio.us/}

%%%%%%%%%%%%%%%%
% NEW CHAPTER! %
%%%%%%%%%%%%%%%%
\chapter{Introdução}

\begin{chapquote}{Autor desconhecido}
``Se você der um programa a alguém, você irá
frustrá-lo por um dia; se você ensiná-lo a
programar, você o frustrará por toda a vida.''

\end{chapquote}
\label{cap1}
O objetivo deste texto é fornecer uma referência para cursos de linguagem assembly e programação de sistemas de nível universitário. Especificamente, este texto aborda o conjunto de instruções x86-64\footnote{Para obter mais informações, consulte: \url{http://pt.wikipedia.org/wiki/X86-64}} para a popular classe de processadores x86-64 usando o sistema operacional (SO) Ubuntu de 64 bits. Embora o código fornecido e vários exemplos devam funcionar em qualquer sistema operacional baseado em Linux de 64 bits, eles foram testados apenas no Ubuntu 14.04 LTS (64 bits).

O x86-64 é um projeto de um Computador com um Conjunto Complexo de Instruções (ou, do inglês: Complex Instruction Set Computing - CISC\footnote{Para obter mais informações, consulte: https://pt.wikipedia.org/wiki/CISC}). Isso se refere à filosofia de design do processador interno. Os processadores CISC normalmente incluem uma ampla variedade de instruções (às vezes sobrepostas), tamanhos variados de instruções e uma ampla gama de modos de endereçamento. O termo foi cunhado retroativamente em contraste com Computador com um conjunto reduzido de instruções (ou, do inglês: Reduced Instruction Set Computer (RISC\footnote{Para obter mais informações, consulte: https://pt.wikipedia.org/wiki/RISC}).

\section{Pré-requisitos}
Deve-se notar que o texto não visa aprender a programar. Supõe-se que o leitor já se tornou proficiente em uma linguagem de programação de alto nível. Especificamente, o texto é geralmente voltado para uma linguagem de alto nível baseada em C compilada, como C, C ++ ou Java. Muitas das explicações e exemplos assumem que o leitor já está familiarizado com conceitos de programação, como declarações, operações aritméticas, estruturas de controle, iteração, chamadas de função, funções, indireção (ou seja, ponteiros) e questões de escopo variável.

Além disso, o leitor deve se sentir confortável ao usar um sistema operacional baseado em Linux, incluindo o uso da linha de comando. Se o leitor for novo no Linux, a seção Referências adicionais contém links para algumas documentações úteis.

\section{O que é a linguagem assembly}
A pergunta típica feita pelos alunos é 'por que aprender \textit{assembly}?'. Antes de abordar essa questão, vamos esclarecer o que exatamente é a linguagem \textit{assembly}.

A linguagem \textit{assembly} é específica da máquina. Por exemplo, o código escrito para um processador x86-64 não será executado em um processador diferente, como um processador RISC (popular em tablets e smartphones).

A linguagem \textit{assembly} é uma linguagem de ``baixo nível'' e fornece a interface de instrução básica para o processador do computador. A linguagem \textit{assembly} é o mais próximo do processador que você pode chegar como um programador. Os programas escritos em uma linguagem de alto nível são traduzidos para a linguagem \textit{assembly} para que o processador execute o programa. A linguagem de alto nível é uma abstração entre a linguagem e as instruções reais do processador. Como tal, a ideia de que ``o assembly está morto`'' é um absurdo.

A linguagem \textit{assembly} fornece controle direto dos recursos do sistema. Isso envolve definir os registros do processador, acessar locais de memória e fazer interface com outros elementos de hardware. Isso requer uma compreensão significativamente mais profunda de exatamente como o processador e a memória funcionam.

\section{Por que aprender a linguagem assembly?}
O objetivo deste texto é fornecer uma introdução abrangente à programação em linguagem \textit{assembly}. As razões para aprender a linguagem \textit{assembly} são mais sobre como entender como um computador funciona, em vez de desenvolver programas grandes. Já que a linguagem \textit{assembly} é específica da máquina, a falta de portabilidade é muito limitante para projetos de programação.

O processo de realmente aprender a linguagem \textit{assembly} envolve escrever programas não triviais para realizar ações específicas de baixo nível, incluindo operações aritméticas, chamadas de função, usando variáveis locais dinâmicas na pilha e interação do sistema operacional para atividades como entrada/saída. Apenas olhar para pequenos programas em linguagem \textit{assembly} não será suficiente.

No longo prazo, aprender os princípios básicos, incluindo a linguagem \textit{assembly}, é o que faz a diferença entre um técnico de codificação incapaz de lidar com linguagens em mudança e um cientista da computação que é capaz de se adaptar às tecnologias em constante mudança.

As seções a seguir fornecem alguns detalhes sobre os vários motivos mais específicos para aprender a linguagem \textit{assembly}.

\subsection{Obter uma melhor compreensão dos problemas de arquitetura}
Aprender e passar algum tempo trabalhando no nível da linguagem assembly fornece uma compreensão mais rica da arquitetura de computador subjacente. Isso inclui o conjunto de instruções básicas, registros do processador, endereçamento de memória, interface de hardware e entrada/saída. Uma vez que, em última análise, todos os programas são executados nesse nível, conhecer os recursos da linguagem \textit{assembly} fornece \textit{insights} úteis sobre o que é possível, o que é fácil e o que pode ser mais difícil ou mais lento.

\subsection{Compreender a cadeia de ferramentas}
A cadeia de ferramentas é o nome do processo de pegar o código escrito por um humano e convertê-lo em algo que o computador possa executar diretamente. Isso inclui o compilador, ou \textit{assembler} (montador) em nosso caso, o \textit{linker} (vinculador), o \textit{loader} (carregador) e o \textit{debugger} (depurador). Em referência à compilação, os programadores iniciantes são informados ``apenas faça isso`'', com pouca explicação sobre a complexidade envolvida no processo. Trabalhar em um nível inferior pode ajudar a fornecer a base para a compreensão e apreciação dos detalhes da cadeia de ferramentas.

\subsection{Melhorar as habilidades de desenvolvimento de algoritmo}
Trabalhar com linguagem \textit{assembly} e escrever programas de baixo nível ajuda os programadores a melhorar as habilidades de desenvolvimento de algoritmos, praticando com uma linguagem que requer mais reflexão e mais atenção aos detalhes. No caso altamente improvável de um programa não funcionar na primeira vez, a depuração de linguagem \textit{assembly} também fornece depuração prática e requer uma abordagem mais sutil, pois apenas adicionar um punhado de instruções de saída é mais difícil no nível da linguagem \textit{assembly}. Isso normalmente envolve um uso mais abrangente de um depurador, que é uma habilidade útil para qualquer programador.

\subsection{Melhorar a compreensão das funções/procedimentos}
Trabalhar com a linguagem \textit{assembly} fornece uma compreensão muito melhor de como funcionam as chamadas de função/procedimento. Isso inclui o conteúdo e a estrutura do \textit{frame} de chamada de função, também conhecido como registro de ativação. Dependendo da instância específica, o registro de ativação pode incluir argumentos baseados em pilha, registros preservados e/ou variáveis locais dinâmicas de pilha. Existem algumas implicações significativas de implementação e segurança em relação às variáveis locais dinâmicas da pilha que são mais bem compreendidas trabalhando em um nível baixo. Devido às implicações de segurança, seria apropriado lembrar aos leitores de sempre usarem seus poderes para o bem. Além disso, o uso da pilha e do \textit{frame} de chamada associado é a base para a recursão e a compreensão da implementação bastante direta das funções recursivas.

\subsection{Obter uma compreensão do buffer de E/S}
Em uma linguagem de alto nível, as instruções de entrada/saída e as operações de \textit{buffer} associadas podem parecer mágicas. Trabalhar no nível da linguagem \textit{assembly} e realizar algumas operações de entrada/saída de baixo nível fornece uma compreensão mais detalhada de como a entrada/saída e o armazenamento em \textit{buffer} realmente funcionam. Isso inclui as diferenças entre entrada/saída interativa, entrada/saída de arquivo e os serviços de sistema operacional associados.

\subsection{Compreender o escopo do compilador}
Programar em linguagem \textit{assembly}, depois de já ter aprendido uma linguagem de alto nível, ajuda a garantir que os programadores entendam o escopo e as capacidades de um compilador. Especificamente, isso significa aprender o que o compilador faz e o que não faz em relação à arquitetura do computador.

\subsection{Introdução conceitos de multiprocessamento}
Este texto também fornecerá uma breve introdução aos conceitos de multiprocessamento. Os conceitos gerais de programação distribuída e multi-core são apresentados com o foco sendo colocado na memória compartilhada, processamento encadeado. É crença do autor que a verdadeira compreensão das questões sutis associadas ao encadeamento, como memória compartilhada e condições de corrida, é mais facilmente compreendida em baixo nível.

\subsection{Introdução aos conceitos de processamento de interrupção}
O mecanismo fundamental subjacente no qual os computadores multiusuário modernos funcionam é baseado em interrupções. Trabalhar em um nível inferior é o melhor lugar para fornecer uma introdução aos conceitos básicos associados ao tratamento de interrupções, identificadores de serviço de interrupção e interrupções de vetor.

\section{Referências Adicionais}
Algumas referências importantes para informações adicionais são observadas nas seções a seguir. Essas referências fornecem informações muito mais extensas e detalhadas.

Se algum desses locais mudar, uma pesquisa na web poderá encontrar o novo local.

\subsection{Referências para o Ubuntu}
Há documentação significativa disponível para o sistema operacional Ubuntu. O principal guia do usuário é o seguinte:
\begin{itemize}
	\item \href{https://ubuntuforum-br.org/}{Fórum Ubuntu}
	\item \href{https://help.ubuntu.com/stable/ubuntu-help/getting-started.html}{Primeiros Passos}
\end{itemize}

Além disso, existem muitos outros sites dedicados a fornecer ajuda para usar o Ubuntu (ou outros sistemas operacionais baseados em Linux).

\subsection{Referências de linha de comando BASH}
\textbf{BASH} é o shell padrão do Ubuntu. O leitor deve estar familiarizado com as operações básicas de linha de comando. Algumas referências adicionais são as seguintes:
\begin{itemize}
	\item \href{https://www.hostinger.com.br/tutoriais/comandos-linux}{35 Comandos Linux Que Todo Usuário Deve Conhecer}
	\item \href{http://leg.ufpr.br/~fernandomayer/aulas/ce083/shell-linux.html}{Introdução ao Shell Linux}
\end{itemize}

Além disso, existem muitos outros sites dedicados a fornecer informações sobre o shell de comando \textbf{BASH}.

\subsection{Referências de arquitetura}
Algumas referências importantes publicadas pela Intel fornecem uma descrição técnica detalhada da arquitetura e do ambiente de programação dos processadores Intel que oferecem suporte às arquiteturas IA-32 e Intel 64.
\begin{itemize}
	\item \href{https://www.intel.com/content/dam/www/public/us/en/documents/manuals/64-ia-32-architectures-software-developer-vol-1-manual.pdf}{Intel\textregistered\space 64 and IA-32 Architectures	Software Developer’s Manual: Basic
		Architecture.}
	\item \href{http://www.intel.com/content/dam/www/public/us/en/documents/manuals/64-ia-32-architectures-software-developer-instruction-set-reference-manual-325383.pdf}{Intel\textregistered\space 64 and IA-32 Architectures Software Developer’s Manual: Instruction
		Set Reference.}
	\item  \href{https://www.intel.com/content/dam/www/public/us/en/documents/manuals/64-ia-32-architectures-software-developer-system-programming-manual-325384.pdf}{Intel\textregistered\space 64 and IA-32 Architectures Software Developer’s Manual: System
		Programming Guide.}
\end{itemize}

Se os links incorporados não funcionarem, uma pesquisa na Internet pode ajudar a encontrar o novo local.

\subsection{Referências de cadeia de ferramentas}
A cadeia de ferramentas inclui o montador, o vinculador, o carregador e o depurador. O Capítulo \ref*{Ferramentasdedesenvolvimento}, Cadeia de ferramentas, fornece uma visão geral da cadeia de ferramentas usada neste texto. As referências a seguir fornecem informações e documentação mais detalhadas.

\subsubsection{Referências YASM}
O assembler YASM é um assembler de código aberto comumente disponível em sistemas baseados em Linux. As referências \textbf{YASM} são as seguintes:
\begin{itemize}
	\item \href{http://yasm.tortall.net/}{Site Web Yasm}
	\item \href{http://yasm.tortall.net/Guide.html}{Documentação Yasm}
\end{itemize}
Informações adicionais sobre YASM podem estar disponíveis em vários sites de linguagem assembly e podem ser encontradas por meio de uma pesquisa na Internet.

\subsubsection{Referências do Debugger DDD}
O depurador DDD é um depurador de código aberto capaz de oferecer suporte a linguagem assembly.
\begin{itemize}
	\item \href{http://www.gnu.org/software/ddd/}{Site Web DDD}
	\item \href{http://www.gnu.org/software/ddd/manual/}{Documentação DDD}
\end{itemize}
Informações adicionais sobre DDD podem estar em vários sites de linguagem assembly e podem ser encontradas em uma pesquisa na Internet.


\chapter{Visão geral da arquitetura}

\begin{chapquote}{Autor desconhecido}
``Aviso, teclado não encontrado. Pressione Enter para continuar.''
\end{chapquote}

Este capítulo apresenta uma visão geral básica da arquitetura \textbf{x86-64}. Para uma explicação mais detalhada, consulte as referências adicionais observadas no Capítulo \ref{cap1}, Introdução.

\section{Visão geral da arquitetura}
Os componentes básicos de um computador incluem uma unidade de processamento central (CPU), armazenamento primário ou memória de acesso aleatório (RAM), armazenamento secundário, dispositivos de entrada/saída (por exemplo, tela, teclado, mouse) e uma interconexão conhecida como barramento.

Um diagrama básico da arquitetura do computador é a da Figura \ref{arquiteturacomputador}:
\begin{figure}[ht]
	\includegraphics[width=\linewidth]{imagens/arquitetura}
	\caption{Arquitetura de um computador}
	\label{arquiteturacomputador}
\end{figure}

A arquitetura é normalmente conhecida como Arquitetura Von Neumann\footnote{Para obter mais informações, consulte: https://pt.wikipedia.org/wiki/Arquitetura\_de\_von\_Neumann}, ou arquitetura Princeton, e foi descrita em 1945 pelo matemático e físico John von Neumann.

Programas e dados são normalmente armazenados em armazenamento secundário (por exemplo, unidade de disco ou unidade de estado sólido). Quando um programa é executado, ele deve ser copiado do armazenamento secundário para o armazenamento primário ou memória principal (RAM). A CPU executa o programa a partir do armazenamento primário ou RAM.

O armazenamento primário ou memória principal também é conhecido como \textit{memória volátil}, pois quando a energia é removida, a informação não é retida e, portanto, perdida. O armazenamento secundário é conhecido como \textit{memória não volátil}, uma vez que as informações são retidas quando desligadas.

Por exemplo, considere armazenar um trabalho de conclusão de curso em armazenamento secundário (ou seja, disco). Quando o usuário começa a escrever ou editar o trabalho de conclusão de curso, ele é copiado do meio de armazenamento secundário para o armazenamento primário (isto é, RAM ou memória principal). Quando concluído, a versão atualizada normalmente é armazenada de volta no armazenamento secundário (ou seja, disco). Se você já perdeu energia durante a edição de um documento (presumindo que não há bateria ou fonte de alimentação ininterrupta), perder o trabalho não salvo certamente esclarecerá a diferença entre memória volátil e não volátil.

\section{Tamanhos de armazenamento de dados}
A arquitetura x86-64 oferece suporte a um conjunto específico de elementos de tamanho de armazenamento de dados, todos baseados em potências de dois. Os tamanhos de armazenamento suportados são os da Tabela \ref{tab:tamanhosdearmazenamento}:

\begin{table}[h]
	\centering
	\begin{tabular}{|l|c|c|}
		\hline
		\rowcolor[HTML]{C0C0C0} 
		{\color[HTML]{000000} } Armazenamento & {\color[HTML]{000000} Tamanho (bits)} & {\color[HTML]{000000} Tamanho (bytes)} \\ \hline
		Byte& 8 bits & 1 byte\\ \hline
		Word & 16 bits & 2 bytes\\ \hline
		Double-word & 32 bits & 4 bytes\\ \hline
		Quadword & 64 bits & 8 bytes\\ \hline
		Double quadword & 128 bits & 16 bytes  \\ \hline
	\end{tabular}
	\caption{Tamanhos de armazenamento.}
	\label{tab:tamanhosdearmazenamento}
\end{table}

Listas ou arrays (conjuntos de memória) podem ser reservados em qualquer um desses tipos.

Esses tamanhos de armazenamento têm uma correlação direta com as declarações de variáveis em linguagens de alto nível (por exemplo, C, C ++, Java, etc.).

Por exemplo, as declarações C/C ++ são mapeadas da seguinte maneira:
\begin{table}[h]
	\centering
	\begin{tabular}{|l|l|c|c|}
		\hline
		\rowcolor[HTML]{C0C0C0} 
		{\color[HTML]{000000} } Declaração em C/C++ & {\color[HTML]{000000} Armazenamento} & {\color[HTML]{000000} Tamanho (bits)} & {\color[HTML]{000000} Tamanho (bytes)} \\ \hline
		char & Byte& 8 bits & 1 byte\\ \hline
		short & Word & 16 bits & 2 bytes\\ \hline
		int & Double-word & 32 bits & 4 bytes\\ \hline
		unsigned int & Double-word & 32 bits & 4 bytes\\ \hline
		long\footnotemark & Quadword & 64 bits & 8 bytes\\ \hline
		long long  & Quadword & 64 bits & 8 bytes\\ \hline
		char *  & Quadword & 64 bits & 8 bytes\\ \hline
		int * & Quadword & 64 bits & 8 bytes\\ \hline
		float & Double quadword & 128 bits & 16 bytes  \\ \hline
		double & Double quadword & 128 bits & 16 bytes  \\ \hline
	\end{tabular}

	\caption{Declarações em C/C++.}
	\label{tab:declaracoesemC}
\end{table}
\footnotetext[2]{Observe que a declaração do tipo 'long' depende do compilador. O tipo mostrado é para compiladores gcc e g ++.}

O asterisco indica uma variável de endereço. Por exemplo, \textbf{int *} significa o endereço de um inteiro. Outras linguagens de alto nível geralmente têm mapeamentos semelhantes.

\section{Unidade central de processamento}
A \textbf{Unidade de Processamento Central}\footnote{Para obter mais informações, consulte: https://pt.wikipedia.org/wiki/Unidade\_central\_de\_processamento} (CPU) é normalmente referida como o ``cérebro'' do computador, uma vez que é onde os cálculos reais são realizados. A CPU está alojada em um único chip, às vezes chamado de processador, chip ou \textit{die}\footnote{Para obter mais informações, consulte: https://pt.wikipedia.org/wiki/Die\_(circuito\_integrado)}. A imagem da capa mostra tal CPU.

O chip da CPU inclui várias unidades funcionais, incluindo a Unidade Lógica e Aritmética\footnote{Para obter mais informações, consulte: https://pt.wikipedia.org/wiki/Unidade\_lógica\_e\_aritmética} (ULA), que é a parte do chip que realmente executa os cálculos aritméticos e lógicos. Para dar suporte à ULA, os registros do processador\footnote{Para obter mais informações, consulte: https://pt.wikipedia.org/wiki/Registrador\_(informática)} e a memória cache\footnote{Para obter mais informações, consulte: https://pt.wikipedia.org/wiki/Cache} também estão incluídos `no dado''` (termo para dentro do chip). Os registros da CPU e a memória cache são descritos nas seções subsequentes.

\subsection{Registradores da CPU}
Um registrador da CPU, ou apenas registrador, é um armazenamento temporário ou local de trabalho embutido na própria CPU (separado da memória). Os cálculos são normalmente executados pela CPU usando registradores.

\subsubsection{Registradores de uso geral (GPRs)}
Existem dezesseis registradores de uso geral (General Purpose Registers - GPRs) de 64 bits. Os GPRs são descritos na tabela a seguir. Um registro GPR pode ser acessado com todos os 64 bits ou alguma parte ou subconjunto acessado (vide Tabela \ref{tab:registradores}).

\begin{table}[h]
	\centering
	\begin{tabular}{|c|c|c|c|}
		\hline
		\rowcolor[HTML]{C0C0C0} 
		Registrador & 32 bits & 16 bits & 8 bits \\ 
		\rowcolor[HTML]{C0C0C0} 
		de 64 bits & mais baixos & mais baixos & mais baixos \\ \hline
		rax & eax& ax & al\\ \hline
		rbx & ebx& bx & bl\\ \hline
		rcx & ecx& cx & cl\\ \hline
		rdx & edx& dx & dl\\ \hline
		rsi & esi& si & sil\\ \hline
		rdi & edi& di & dil\\ \hline
		rbp & ebp& bp & bpl\\ \hline
		rsp & esp& sp & spl\\ \hline
		r8 & r8d& r8w & r8b\\ \hline
		r9 & r9d& r9w & r9b\\ \hline
		r10 & r10d& r10w & r10b\\ \hline
		r11 & r11d& r11w & r11b\\ \hline
		r12 & r12d& r12w & r12b\\ \hline
		r13 & r13d& r13w & r13b\\ \hline
		r14 & r14d& r14w & r14b\\ \hline
		r15 & r15d& r15w & r15b\\ \hline
	\end{tabular}
	
	\caption{Registradores da CPU.}
	\label{tab:registradores}
\end{table}

Além disso, alguns dos registradores GPR são usados para fins específicos, conforme descrito nas seções posteriores.

Ao usar tamanhos de elementos de dados menores que 64 bits (ou seja, 32 bits, 16 bits ou 8 bits), a parte inferior do registro pode ser acessada usando um nome de registro diferente, conforme mostrado na tabela.

Por exemplo, ao acessar as partes inferiores do registrador \textbf{rax} de 64 bits, o layout é o seguinte:

\begin{figure}[ht]
	\includegraphics[width=0.8\linewidth]{imagens/layoutrax}
	\caption{Layout do registrador rax.}
	\label{layoutrax}
\end{figure}

Conforme mostrado no diagrama, os primeiros quatro registros, \textbf{rax}, \textbf{rbx}, \textbf{rcx} e \textbf{rdx} também permitem que os bits 8-15 sejam acessados com os nomes de registro \textbf{ah}, \textbf{bh}, \textbf{ch} e \textbf{dh}. Com exceção de \textbf{ah}, eles são fornecidos para suporte de legado e não serão usados neste texto.

A capacidade de acessar partes do registro significa que, se o registro quadword \textbf{rax}  for definido como $ 50.000.000.000_{10} $ (cinquenta bilhões), o registro \textbf{rax} conterá o seguinte valor em \textit{hex}.

\begin{center}
	\textbf{rax = 0000 000B A43B 7400}
\end{center}

Se uma operação subsequente definir o registro da palavra \textbf{ax} para $ 50.000_{10} $ (cinquenta mil, que é $ C350_{16} $), o registro \textbf{rax} conterá o seguinte valor em \textit{hex}.

\begin{center}
	\textbf{rax = 0000 000B A43B C350}
\end{center}

Nesse caso, quando a parte inferior \textbf{ax} de 16 bits do registrador \textbf{rax} de 64 bits é definida, os 48 bits superiores não são afetados. Observe a mudança no \textbf{AX} (de $ 7400_{16} $ para $ C350_{16} $).

Se uma operação subsequente definir o registro \textbf{al} de tamanho de byte para $ 50_{10} $ (cinquenta, que é $ 32_{16} $), o registro \textbf{rax} conterá o seguinte valor em hexadecimal.

\begin{center}
	\textbf{rax = 0000 000B A43B C332}
\end{center}

Quando a parte inferior \textbf{al} de 8 bits do registro \textbf{rax} de 64 bits é definida, os 56 bits superiores não são afetados. Observe a mudança em \textbf{AL} (de $ 50_{16} $ para $ 32_{16}$).

Para operações de registrador de 32 bits, os 32 bits superiores são apagados (definidos para zero). Geralmente, isso não é um problema, pois as operações em registradores de 32 bits não usam os 32 bits superiores do registrador. Para valores não sinalizados, isso pode ser útil para converter de 32 bits para 64 bits.
No entanto, isso não funcionará para conversões sinalizadas de valores de 32 bits para 64 bits.
Especificamente, ele fornecerá resultados potencialmente incorretos para valores negativos de números.
Consulte o Capítulo \ref{cap:representacaodedados}, \textit{Representação de dados}, para obter informações adicionais sobre a representação de valores com sinais.

\subsubsection{Stack Pointer Register (RSP)}
Um dos registradores da CPU, \textbf{rsp}, é usado para apontar para o topo atual da pilha. O registrador \textbf{rsp} não deve ser usado para dados ou outros usos. Informações adicionais sobre as operações de pilha e pilha são fornecidas no Capítulo \ref{cap:pilhadeprocessos}, Pilha de processos.

\subsubsection{Base Pointer Register (RBP)}
Um dos registradores da CPU, \textbf{rbp}, é usado como um ponteiro base durante as chamadas de função. O registro \textbf{rbp} não deve ser usado para dados ou outros usos. Informações adicionais sobre as funções e chamadas de função são fornecidas no Capítulo \ref{cap:Funções}, Funções.


\subsubsection{Instruction Pointer Register (RIP)}
Além dos GPRs, existe um registro especial, \textbf{rip}, que é usado pela CPU para apontar a próxima instrução a ser executada. Especificamente, uma vez que o \textbf{rip} aponta para a próxima instrução, isso significa que a instrução apontada por \textbf{rip} e mostrada no depurador ainda não foi executada. Esta é uma distinção importante que pode ser confusa ao revisar o código em um depurador.

\subsubsection{Flag Register (rFlags)}
O registrador de sinalização, \textbf{rFlags}, é usado para informações de status e controle da CPU. O registro \textbf{rFlag} é atualizado pela CPU após cada instrução e não pode ser acessado diretamente pelos programas. Este registro armazena informações de status sobre a instrução que acabou de ser executada. Dos 64 bits do registrador \textbf{rFlag}, muitos são reservados para uso futuro.

A Tabela \ref{tab:flags} mostra alguns dos bits de status no registrador de sinalizador.

\begin{table}[h]
	\centering
	\begin{tabular}{|c|c|c|l|}
		\hline
		\rowcolor[HTML]{C0C0C0} 
		\textbf{Nome} & \textbf{Símbolo} & \textbf{Bit} & \textbf{Uso} \\ 
		Carry & CF & 0  & Usado para indicar se a operação anterior\\
		& & & resultou em um transporte (vai um).\\ \hline
		Paridade & PF& 2 & Usado para indicar se o último byte tem\\
		& & &  um número par de 1 (ou seja, paridade par).\\ \hline
		Ajuste & AF& 4 & Usado para suportar operações decimais\\
		& & &  codificadas binárias.\\ \hline
		Zero & ZF& 6 & Usado para indicar se a operação anterior\\
		& & &  resultou em um resultado zero.\\ \hline
		Sinal & SF& 7 & Utilizado para indicar se o resultado\\
		& & &  da operação anterior resultou em 1 no bit\\
		& & &  mais significativo (indicando negativo no\\
		& & &  contexto de dados sinalizados).\\ \hline
		Direção & DF& 10 & Usado para especificar a direção (incremento\\
		& & &  ou decremento) para algumas operações de\\
		& & &  string.\\ \hline
		Overflow & OF& 11 & Usado para indicar se a operação anterior\\
		& & &  resultou em um estouro.\\ \hline
	\end{tabular}
	
	\caption{Flags.}
	\label{tab:flags}
\end{table}

Existem vários bits adicionais não especificados neste texto. Mais informações podem ser obtidas nas referências adicionais observadas no Capítulo \ref{cap1}, \textit{Introdução}.

\subsubsection{Registradores XMM}
Há um conjunto de registradores dedicados usados para oferecer suporte a operações de ponto flutuante de 64 e 32 bits e instruções de Single Instruction Multiple Data (SIMD). As instruções \textbf{SIMD} permitem que uma única instrução seja aplicada simultaneamente a vários itens de dados. Usado de forma eficaz, isso pode resultar em um aumento significativo de desempenho. As aplicações típicas incluem algum processamento gráfico e processamento de sinal digital.

Os registradores \textbf{XMM} são apresentados na Tabela \ref{tab:xmm}.


\begin{table}[h]
	\centering
	\begin{tabular}{|c|}
		\hline
		\rowcolor[HTML]{C0C0C0} 
		\textbf{Registradores de 128 bits} \\ \hline
		xmm0 \\ \hline
		xmm1 \\ \hline
		xmm2 \\ \hline
		xmm3 \\ \hline
		xmm4 \\ \hline
		xmm5 \\ \hline
		xmm6 \\ \hline
		xmm7 \\ \hline
		xmm8 \\ \hline
		xmm9 \\ \hline
		xmm10 \\ \hline
		xmm11 \\ \hline
		xmm12 \\ \hline
		xmm13 \\ \hline
		xmm14 \\ \hline
		xmm15 \\ \hline
	\end{tabular}
	
	\caption{Registradores xmm.}
	\label{tab:xmm}
\end{table}

Observe que alguns dos processadores X86-64 mais recentes suportam registros \textbf{XMM} de 256 bits. Isso não será uma questão para os programas neste texto.

Além disso, os registradores \textbf{XMM} são usados para oferecer suporte a \textit{Streaming SIMD Extensions} (SSE). As instruções \textbf{SSE} estão fora do escopo deste texto. Mais informações podem ser obtidas nas referências da Intel (conforme observado no Capítulo \ref{cap1}, \textit{Introdução}).

\subsection{Memória cache}
A memória cache é um pequeno subconjunto do armazenamento primário ou RAM localizado no chip da CPU. Se um local de memória for acessado, uma cópia do valor é colocada no cache. Os acessos subsequentes a esse local da memória que ocorrem em rápida sucessão são recuperados do local do cache (interno ao chip da CPU). Uma leitura de memória envolve o envio do endereço por meio do barramento para o controlador de memória, que obterá o valor no local de memória solicitado e o enviará de volta pelo barramento. Comparativamente, se um valor estiver no cache, será muito mais rápido acessar esse valor.

Um \textbf{acerto} de cache ocorre quando os dados solicitados podem ser encontrados em um cache, enquanto uma \textbf{falha} de cache ocorre quando não pode. As ocorrências de cache são atendidas pela leitura de dados do cache, o que é mais rápido do que ler da memória principal. Quanto mais solicitações podem ser atendidas do cache, mais rápido o sistema normalmente executará. Gerações sucessivas de chips de CPU aumentaram a memória cache e aprimoraram as estratégias de mapeamento de cache para melhorar o desempenho geral.

Um diagrama de blocos de uma configuração típica de chip de CPU é o dada na Figura \ref{fig:cache} :
\begin{figure}[ht]
\begin{center}
		\includegraphics[width=0.8\linewidth]{imagens/cache}
\end{center}
	\caption{Diagrama em blocos de memória cache}
	\label{fig:cache}
\end{figure}

Os designs de chip atuais normalmente incluem um cache \textbf{L1} por núcleo e um cache \textbf{L2} compartilhado. Muitos dos chips de CPU mais recentes terão um cache \textbf{L3} adicional.

Como pode ser observado no diagrama, todos os acessos à memória passam por cada nível de cache. Como tal, há um potencial para cópias múltiplas e duplicadas do valor (registro de CPU, cache L1, cache L2 e memória principal). Essa complicação é gerenciada pela CPU e não é algo que o programador possa alterar. Compreender o cache e o ganho de desempenho associado é útil para entender como um computador funciona.

\section{Memória principal}
A memória pode ser vista como uma série de bytes, um após o outro. Ou seja, a memória é endereçável por byte. Isso significa que cada endereço de memória contém um byte de informação. Para armazenar uma palavra dupla, são necessários quatro bytes que usam quatro endereços de memória.

Além disso, a arquitetura é \textbf{little-endian}. Isso significa que o Byte menos significativo (\textit{Least Significant Byte}, LSB)) é armazenado no endereço de memória mais baixo. O Byte mais significativo (\textit{Most Significant Byte}, MSB) é armazenado no local de memória mais alto.

Para uma palavra dupla (32 bits), o MSB e o LSB são alocados conforme mostrado na Figura \ref{fig:msb}.
\begin{figure}[ht]
	\begin{center}
		\includegraphics[width=\linewidth]{imagens/msb}
	\end{center}
	\caption{MSB e LSB}
	\label{fig:msb}
\end{figure}

Por exemplo, assumindo o valor de $ 5.000.000_{10} $ ($ 004C4B40_{16} $), deve ser colocado em uma variável de palavra dupla chamada \textbf{var1}.

\tikzstyle{memoria} = [rectangle, draw, fill=blue!20, 
text width=4em, text centered]

Para uma arquitetura \textit{little-endian}, a imagem da memória seria a da Figura \ref{fig:littleendien}:
\begin{figure}[ht]
	\begin{center}
	\begin{tikzpicture}
\node [text width=5em,xshift=0.5cm,] (nome) {\textbf{Nome da variável}};
\node [text width=5em, xshift=3cm, yshift=-0.2cm] (valor)  {\textbf{Valor}};
\node [text width=5em, xshift=4.5cm, yshift=-0.2cm] (endereco) {\textbf{Endereço}};
\node [memoria, xshift=2.5cm, yshift=-0.8cm ]  {00};
\node [xshift=4.5cm, yshift=-0.8cm ]  {0x00000010};
\node [memoria, xshift=2.5cm, yshift=-1.3cm ]  {89};
\node [xshift=4.5cm, yshift=-1.3cm ]  {0x0000000F};
\node [memoria, xshift=2.5cm, yshift=-1.8cm ]  {54};
\node [xshift=4.5cm, yshift=-1.8cm ]  {0x0000000E};
\node [memoria, xshift=2.5cm, yshift=-2.3cm ]  {40};
\node [xshift=0.5cm, yshift=-2.3cm ]  {var1 $\Rightarrow$};
\node [xshift=4.5cm, yshift=-2.3cm ]  {0x0000000D};
\node [memoria, xshift=2.5cm, yshift=-2.8cm ]  {?};
\node [xshift=4.5cm, yshift=-2.8cm ]  {0x0000000C};
\end{tikzpicture}
	\end{center}
	\caption{Layout de dados Little-Endien}
	\label{fig:littleendien}
\end{figure}

Com base na arquitetura \textit{little-endian}, o \textbf{LSB} é armazenado no endereço de memória mais baixo e o \textbf{MSB} é armazenado no local de memória mais alto.

\section{Layout de memória}
O layout geral da memória para um programa é mostrado na Figura \ref{fig:memlayout}.
\begin{figure}[ht]
	\begin{center}
		\includegraphics[width=0.8\linewidth]{imagens/memlayout}
	\end{center}
	\caption{Layout geral de memória}
	\label{fig:memlayout}
\end{figure}

A seção reservada não está disponível para programas do usuário. A seção de texto (ou código) é onde a linguagem de máquina\footnote{Para obter mais informações, consulte: https://pt.wikipedia.org/wiki/Código\_de\_máquina} (ou seja, os \textbf{1s} e \textbf{0s} que representam o código) é armazenada. A seção de dados é onde os dados inicializados são armazenados. Isso inclui variáveis declaradas que receberam um valor inicial no momento da montagem. A seção de dados não inicializada, normalmente chamada de seção BSS, é onde as variáveis declaradas que não receberam um valor inicial são armazenadas. Se acessado antes de ser definido, o valor não será significativo. O \textit{heap} é onde os dados alocados dinamicamente serão armazenados (se solicitados). A pilha (stack) começa com muita memória e cresce para baixo.

As seções posteriores fornecerão detalhes adicionais para as seções de texto e dados.


\section{Hierarquia de Memória}
A fim de compreender completamente os vários níveis de memória diferentes e uso associado, é útil revisar a hierarquia de memória\footnote{Para obter mais informações, consulte: \href{https://pt.wikipedia.org/wiki/Hierarquia\_de\_mem\%C3\%B3ria}{Hierarquia de Memória, na Wikipedia}}. Em termos gerais, a memória mais rápida é mais cara e os blocos de memória mais lentos são menos caros. Os registros da CPU são pequenos, rápidos e caros. Os dispositivos de armazenamento secundários, como unidades de disco e unidades de estado sólido (Solid State Devices, SSDs), são maiores, mais lentos e mais baratos. O objetivo geral é equilibrar desempenho e custo.

Uma visão geral da hierarquia da memória é dada na Figura \ref{fig:hierarquia}:
\begin{figure}[ht]
	\begin{center}
		\includegraphics[width=0.8\linewidth]{imagens/hierarquia}
	\end{center}
	\caption{Hierarquia de memória}
	\label{fig:hierarquia}
\end{figure}

Onde o topo do triângulo representa a memória mais rápida, menor e mais cara. À medida que descemos os níveis, a memória se torna mais lenta, maior e menos cara. O objetivo é usar um equilíbrio eficaz entre a memória pequena, rápida e cara e a memória grande, mais lenta e barata.

Algumas características típicas de desempenho e tamanho são as dadas na Tabela \ref{tab:desempenho}:
\begin{table}[h]
	\centering
	\begin{tabular}{|l|l|l|}
		\hline
		\rowcolor[HTML]{C0C0C0} 
		\textbf{Unidade} & \textbf{Exemplo} & \textbf{Velocidade} \\ 
		\rowcolor[HTML]{C0C0C0} 
		\textbf{de Memória} & \textbf{de tamanho} & \textbf{Típica} \\ 
		Registradores & Registradores de  & $  \sim1 $ ns\\ 
		& 16, 64 bits &\\ \hline
		Memória cache & 4 - 8+ Megabytes  & $  \sim5-60 $ ns\\ 
		& (L1 e L2)& \\ \hline
		Armazenamento primário & 2 – 32+ Gigabytes & $  \sim100-150 $ ns\\ 
		 (i.e, memória principal) & &\\ \hline
		Armazenamento secundário  & 500 Gigabytes --  & $  \sim3-15 $ ms\\ 
		(i.e, disco, SSDs, etc.) & 4+ Terabytes& \\\hline
	\end{tabular}
	
	\caption{Desempenho e tamanho de memórias.}
	\label{tab:desempenho}
\end{table}

Com base nesta tabela, um acesso ao armazenamento primário em 100 nanossegundos
 ($ 100 \times 10^{-9} $) 
 é 30.000 vezes mais rápido do que um acesso ao armazenamento secundário, em 3 milissegundos ($3 \times 10^{-3}$).
 
 As velocidades típicas melhoram com o tempo (e já estão desatualizadas). O ponto chave é que a diferença relativa entre cada unidade de memória é significativa. Essa diferença entre as unidades de memória se aplica mesmo quando \textbf{SSDs} mais novos e mais rápidos estão sendo utilizados.

\section{Exercícios}
Abaixo estão algumas perguntas baseadas neste capítulo.

\subsection{Questionário}
\begin{enumerate}
	\item Faça um desenho da Arquitetura de Von Neumann.
	\item  Qual componente de arquitetura conecta a memória à CPU?
	\item  Onde os programas são armazenados quando o computador é desligado?
	\item  Onde os programas devem ser localizados durante a execução?
	\item  Como a memória cache ajuda no desempenho geral?
	\item  Quantos bytes um inteiro C ++ declarado com a declaração \textbf{int} usa?
	\item  Na arquitetura \textit{Intel X86-64}, quantos \textbf{bytes} podem ser armazenados em cada endereço?
	\item Dado o hexadecimal de 32 bits $ 004C4B40_{16} $, qual é:
	\begin{enumerate}
		\item Byte menos significativo (LSB)
		\item Byte mais significativo (MSB)
	\end{enumerate}
    \item Dado o hexadecimal de 32 bits $ 004C4B40_{16} $, mostre o layout de memória \textit{little-endian} mostrando cada byte na memória.
	\item  Faça um desenho do layout do registrador \textbf{rax}.
	\item  Quantos bits cada um dos seguintes representa:
	\begin{enumerate}
		\item al
		\item rcx
		\item bx
		\item edx
		\item r11
		\item r8b
		\item sil
		\item r14w
	\end{enumerate}
    \item Qual registrador aponta para a próxima instrução a ser executada?
    \item Qual registrador aponta para o topo da pilha atual?
    \item Se \textbf{al} é definido como $ 05_{16} $ e \textbf{ax} é definido como $ 0007_{16} $, \textbf{eax} é definido como $ 00000020_{16} $, e \textbf{rax} é definido
    a $ 0000000000000000_{16} $, mostrar o conteúdo final completo do registro rax inteiro.
    \item Se o registro \textbf{rax} for definido como $ 81.985.529.216.486.895_{10} $ ($ 123456789ABCDEF_{16} $), qual é o conteúdo dos seguintes registros em hexadecimal?
    \begin{enumerate}
    	\item al
    	\item ax
    	\item eax
    	\item rax
    \end{enumerate}
\end{enumerate}

\chapter{Representação de dados}
\label{cap:representacaodedados}

\begin{chapquote}{Autor desconhecido}
``Existem 10 tipos de pessoas no mundo;
aquelas que entendem binário e aquelas que
não entendem binário.''
\end{chapquote}

A representação de dados se refere a como as informações são armazenadas no computador. Existe um método específico para armazenar números inteiros que é diferente do armazenamento de valores de ponto flutuante, que é diferente do armazenamento de caracteres. Este capítulo apresenta um breve resumo dos esquemas de representação de inteiro, ponto flutuante e ASCII.

Presume-se que o leitor já esteja familiarizado com os sistemas de numeração binária, decimal e hexadecimal.

Deve-se observar que, se não for especificado, um número está na base 10. Além disso, um número precedido por 0x é um valor hexadecimal. Por exemplo, $ 19 = 19_{10}  = 13_{16} = 0x13$.

\section{Representação de Inteiros}
Representar números inteiros refere-se a como o computador armazena ou representa um número na memória. O computador representa números em binários (1's e 0's). No entanto, o computador tem uma quantidade limitada de espaço que pode ser usada para cada número ou variável.
Isso afeta diretamente o tamanho, ou intervalo, do número que pode ser representado. Por exemplo, um byte (8 bits) pode ser usado para representar $ 2^8 $ ou $ 256 $ números diferentes. Esses 256 números diferentes podem não ter sinal (todos positivos), caso em que podemos representar qualquer número entre 0 e 255 (inclusive). Se escolhermos sinais (valores positivos e negativos), podemos representar qualquer número entre $ -128 $ e $ +127 $ (inclusive).

Se esse intervalo não for grande o suficiente para lidar com os valores pretendidos, um tamanho maior deve ser usado. Por exemplo, uma palavra (16 bits) pode ser usada para representar $ 2^{16} $ ou 65.536 valores diferentes, e uma palavra dupla (32 bits) pode ser usada para representar $ 2^{32} $ ou 4.294.967.296 números diferentes. Portanto, se você quiser armazenar um valor de 100.000, será necessária uma palavra dupla.

Como você deve se lembrar de C, C++ ou Java, uma declaração de número inteiro (por exemplo, \textbf{int} \textbf{<variável>}) é uma única palavra dupla que pode ser usada para representar valores entre
 $ -2^{31} $ 
 ($-2.147.483.648$) e 
 $ +2^{31} - 1 $
  ($+ 2.147.483.647$).

A Tabela \ref{tab:faixas} mostra os intervalos associados aos tamanhos típicos:

\begin{table}[h]
	\centering
	\small
	\begin{tabular}{|l|c|c|c|}
		\hline
		\rowcolor[HTML]{C0C0C0} 
		\textbf{Tamanho} & \textbf{Tamanho} & \textbf{Faixa sem sinal} & \textbf{Faixa sinalizada} \\ \hline
		Bytes (8-bits) & $2^8$ & 0 a 255 & $-128$ a $+127$\\ \hline
		Words (16-bits) & $2^{16}$ & 0 a 65.535 & $-32768$ a $+32767$\\ \hline
		Double-words (32-bits) & $2^{32}$ & 0 a 4.294.967.295 & $-^2.147.483.648$ a\\
		&&& $+2.147.483.647$\\ \hline
		Quadword & $2^{64}$ & 0 a $2^{64}-1$ & $-2^{63}$ a $2^{63}-1$\\ \hline
		Double quadword & $2^{128}$ & 0 a 255 & $-128$ a $+127$\\ \hline
	\end{tabular}
	
	\caption{Faixas de valores para tamanhos típicos.}
	\label{tab:faixas}
\end{table}

Para determinar se um valor pode ser representado, você precisará saber o tamanho do elemento de armazenamento (byte, palavra, palavra dupla, palavra quádrupla, etc.) sendo usado e se os valores são sinalizados ou não.

\begin{itemize}
	\item Para representar valores \textit{não sinalizados} dentro do intervalo de um determinado tamanho de armazenamento, o binário padrão é usado. 
	\item Para representar \textit{valores com sinais} dentro do intervalo, o complemento a dois é usado. Especificamente, o processo de codificação de complemento a dois se aplica aos valores na faixa negativa. Para valores dentro da faixa positiva, o binário padrão é usado.
\end{itemize}

Por exemplo, o intervalo de bytes não sinalizados pode ser representado usando uma linha numérica da seguinte maneira:

\vspace{5mm}
\begin{tikzpicture}
\draw[white, thick, |-] (-2.95,0) -- (0,0)  node [black, below=2pt]
{0};
\draw[blue, thick,|-] (0,0) -- (2.5,0);
\draw[blue, thick,|-|] (2.5,0) -- (5,0) node [black, below=2pt]
{255};
\end{tikzpicture}

Por exemplo, o intervalo de bytes sinalizados também pode ser representado usando uma linha numérica da seguinte forma:

\vspace{5mm}
\begin{tikzpicture}
\draw[red, thick, |-]  (-2.5,0) node [black, below=2pt]
{-128} -- (0,0)  node [black, below=2pt]
{0};
\draw[blue, thick,|-|]  (0,0)  -- (2.5,0) node [black, below=2pt]
{+127};
\draw[white, thick,|-|] (2.5,0) -- (5,0);
\end{tikzpicture}

O mesmo conceito se aplica a meias palavras e palavras que possuem intervalos maiores.


Como os valores sem sinal têm um intervalo diferente, apenas positivo, do que os valores com sinal, há sobreposição entre os valores. Isso pode ser muito confuso ao examinar variáveis na memória (com o depurador).

Por exemplo, quando os valores sem sinal e com sinal estão dentro do intervalo positivo de sobreposição (0 a +127):

\begin{itemize}
	\item Uma representação de byte com sinal de $ 12_{10} $ é $ 0x0C_{16} $
	\item Uma representação de byte sem sinal de $-12_{10} $ também é $ 0x0C_{16} $
\end{itemize}

Quando os valores não sinalizados e sinalizados estão fora do intervalo de sobreposição:
\begin{itemize}
	\item Uma representação de byte sinalizado de $ -15_{10} $ é $ 0xF1_{16} $
	\item Uma representação de byte não sinalizado de $ 241_{10} $ também é $ 0xF1_{`} $
\end{itemize}

Essa sobreposição pode causar confusão, a menos que os tipos de dados sejam definidos de forma clara e correta.

\subsection{Complemento a Dois}
A seguir é descrito como encontrar a representação de complemento a dois para valores negativos (valores não positivos).

Para pegar o complemento a dois de um número:
\begin{enumerate}
	\item pegue o complemento a um (negar)
	\item adicione 1 (em binário)
\end{enumerate}

O mesmo processo é usado para codificar um valor decimal no complemento a dois e do complemento a dois de volta ao decimal. As seções a seguir fornecem alguns exemplos.

\subsubsection{Exemplo de Byte}
Por exemplo, para encontrar a representação com tamanho de byte (8 bits), do complemento a dois de $ -9 $ e $ -12$.

\begin{minipage}{.5\textwidth}
\begin{tabular}{|r|c|}
	\hline
	$ 9 (8+1)= $ & 0000 1001\\ \hline
	Passo 1 & 1111 0110\\ \hline
	Passo 2 & 1111 0111\\ \hline
	-9 (em hexa) & \textbf{F7}\\ \hline
\end{tabular}
	
\end{minipage} 
\begin{minipage}{.5\textwidth}
\begin{tabular}{|r|c|}
	\hline
	$ 12 (8+4)= $ & 0000 1100\\ \hline
	Passo 1 & 1111 0011\\ \hline
	Passo 2 & 1111 0100\\ \hline
	-12 (em hexa) & \textbf{F4}\\ \hline
\end{tabular}
	
\end{minipage}%

\vspace{5mm}

Observe que todos os bits para o tamanho fornecido, byte neste exemplo, devem ser especificados.

\subsubsection{Exemplo de Word}
Para encontrar o tamanho da palavra (16 bits), a representação em complemento a dois de $ -18 $ e $ -40$.

\begin{minipage}{.5\textwidth}
	\begin{tabular}{|r|c|}
		\hline
		$ 18 (6+2)= $ & 0000000000010010\\ \hline
		Passo 1 & 1111111111101101\\ \hline
		Passo 2 & 1111111111101110\\ \hline
		-9 (em hexa) & \textbf{FFEE}\\ \hline
	\end{tabular}
	
\end{minipage} 
\begin{minipage}{.5\textwidth}
	\begin{tabular}{|r|c|}
		\hline
		$ 40 (32+8)= $ & 0000000000101000\\ \hline
		Passo 1 & 1111111111010111\\ \hline
		Passo 2 & 1111111111011000\\ \hline
		-12 (em hexa) & \textbf{FFD8}\\ \hline
	\end{tabular}
	
\end{minipage}%

\vspace{5mm}
Observe que todos os bits para o tamanho fornecido, palavras nesses exemplos, devem ser especificados.

\section{Adição não sinalizada e sinalizada}
Conforme observado anteriormente, as representações não assinadas e assinadas podem fornecer diferentes interpretações para o valor final sendo representado. No entanto, as operações de adição e subtração são as mesmas. Por exemplo:

\begin{minipage}{.5\textwidth}
	\begin{tabular}{|r|c|}
		\hline
		$ 241 $ & 11110001\\ \hline
		$+7$ & 00000111\\ \specialrule{2.5pt}{1pt}{1pt}
		$ 248 $ & 11111000\\ \specialrule{2.5pt}{1pt}{1pt}
		$ 248 $ = & \textbf{F8}\\ \hline
	\end{tabular}
	
\end{minipage} 
\begin{minipage}{.5\textwidth}
	\begin{tabular}{|r|c|}
		\hline
		$ -15 $ & 11110001\\ \hline
		$+7$ & 00000111\\ \specialrule{2.5pt}{1pt}{1pt}
		$ -8 $& 11111000\\ \specialrule{2.5pt}{1pt}{1pt}
		$ -8 = $ & \textbf{F8}\\ \hline
	\end{tabular}
	
\end{minipage}%

\vspace{5mm}

O resultado final de \textbf{0xF8} pode ser interpretado como 248 para representação não sinalizada e $ -8 $ para uma representação sinalizada. Além disso, $ 0xF8_{16} $ é o $^{\circ}$ (símbolo de grau) na tabela ASCII.

Como tal, é muito importante ter uma definição clara dos tamanhos (byte, meia palavra, palavra, etc.) e tipos (com ou sem sinal) de dados para as operações que estão sendo realizadas.

\section{Representação de ponto flutuante}
Os problemas de representação para números de ponto flutuante são mais complexos. Há uma série de representações de ponto flutuante para vários intervalos de valor. Para simplificar, examinaremos principalmente o padrão de ponto flutuante \textbf{IEEE 754} de 32 bits.

\subsection{Representação IEEE de 32 bits}
O padrão de ponto flutuante\textbf{ IEEE 754} de 32 bits é definido da seguinte forma:

\noindent
\begin{tikzpicture}[scale=0.4]
\foreach \x in {31,...,0}
{
	\draw (\x,0) +(-.5,-.5) rectangle ++(.5,.5);
	\draw (31-\x, 0) node{{\tiny \x}};	
}
\draw(0,-1) +(-.5,-.5) rectangle ++(.5,.5);
\draw(1,-1) +(-.5,-.5) rectangle ++(7.5,.5);
\draw(9,-1) +(-.5,-.5) rectangle ++(22.5,.5);
\draw(0, -1) node{{\tiny \textbf{s}}};
\draw(4.5, -1) node{{\tiny \textbf{expoente em excesso}}};
\draw(19, -1) node{{\tiny \textbf{fração}}};

\end{tikzpicture}

Onde \textbf{s} é o sinal (0 $\Rightarrow $ positivo e 1 $ \Rightarrow $  negativo). Mais formalmente, isso pode ser escrito como:
\begin{center}
	$ N = (-1)^S \times 1.F \times 2^{E-127} $
\end{center}

Ao representar valores de ponto flutuante, a primeira etapa é converter o valor de ponto flutuante em binário. A tabela a seguir fornece um breve lembrete de como o binário lida com componentes fracionários:

\vspace{5mm}
\begin{center}
	\begin{tabular}{|c|c|c|c|c|c|c|c|c|c|}
	\hline
	         &$2^3$& $2^2$ & $2^1$ & $2^0$ & & $2^{-1}$  & $2^{-2}$ & $2^{-3}$ & \\ \hline
	$\ldots$ & 8   & 4     & 2     & 1     &.& 1/2       &  1/4     & 1/8 & $\ldots$\\ \hline
	         & 0   & 0     & 0     & 0     &.& 0         &  0     & 0 & \\ \hline
\end{tabular}
\end{center}
\vspace{5mm}

Por exemplo, $ 100,101_2 $ seria $ 4,625_{10} $. Para decimais repetidas, o cálculo do valor binário pode ser demorado. No entanto, há um limite, pois os computadores têm tamanhos de armazenamento finitos (32 bits neste exemplo).

A próxima etapa é mostrar o valor em notação científica normalizada em binário. Isso significa que o número deve ter um único dígito inicial diferente de zero à esquerda da vírgula decimal. Por exemplo, $ 8,125_{10} $ é $ 1000,001_2 $ (ou $ 1000,001_2 \times 2^0$) e em notação científica binária normalizada que seria escrita como $ 1,000001 \times 2^3 $ (já que a vírgula decimal foi movida três casas para a esquerda). Claro, se o número fosse $ 0,125_{10} $, o binário seria $ 0,001_2 $ (ou $ 0,001_2 \times 2^0 $) e a notação científica normalizada seria $ 1,0 \times 2_{-3} $ (já que a vírgula decimal foi movida três casas para a direita). Os números após o 1 inicial, não incluindo o 1 inicial, são armazenados justificados à esquerda na parte da fração da palavra dupla.

A próxima etapa é calcular o expoente em excesso, que é o expoente da notação científica normalizada com mais o excesso. O excesso para o padrão de ponto flutuante de 32 bits \textbf{IEEE 754} é $ 127_{10} $. O resultado deve ser convertido para um byte (8 bits) e armazenado na parte expoente em excesso da palavra.

Observe que a conversão da representação de ponto flutuante de 32 bits \textbf{IEEE 754} para o valor decimal é feita ao contrário, no entanto, o 1 inicial deve ser adicionado de volta (já que não é armazenado na palavra). Além disso, o excesso é subtraído (em vez de adicionada).

\subsubsection{Exemplos de representação IEEE de 32 bits}
Esta seção apresenta vários exemplos de codificação e decodificação de representação de ponto flutuante para referência.

\subsubsection{Exemplo: $-7.75_{10} $ }
Por exemplo, para encontrar a representação de ponto flutuante IEEE 754 de 32 bits para $-7.75_{10} $: 

\begin{tabular}{lll}
	1. Determinar o sinal: & $ -7.75 \Rightarrow 1 $ (posto que negativo)\\
	2. converter para binário: & $ -7.75 = -0111.11_2 $\\
	3. notação científica normalizada: & $ =	1.1111 \times 2^2 $\\
	4. calcular expoente em excesso: & $ 2_{10} + 127_{10} = 129_{10} $\\
	e converter para binário: & $ = 10000001_2 $\\
	5. escrever componentes em binário: & \\
    \multicolumn{2}{l}{\begin{tabular}{ccl}
    		sinal & expoente & mantissa\\
    		1 & 10000001 & 11110000000000000000000\\
    \end{tabular}} \\
    \multicolumn{2}{l}{6. converter para hexadecimal (dividido em grupos de 4)} \\
    \multicolumn{2}{l}{11000000111110000000000000000000}\\
    \multicolumn{2}{l}{\begin{tabular}{*8c}
     1100 & 0000 & 1111 & 1000 & 0000 & 0000 & 0000 & 0000\\
     C    & 0    & F    & 8    & 0    & 0    & 0    & 0\\
    	\end{tabular}
    }\\
    7. resultado final: & $ C0F8 0000_16 $\\
	
\end{tabular}




\chapter{Formato de um programa}
\label{cap4:formato}

\begin{chapquote}{Autor desconhecido}
``Eu adoraria mudar o mundo, mas eles não vão me dar o código fonte.''
\end{chapquote}
Este capítulo resume os requisitos de formatação para programas em linguagem assembly. Os requisitos de formatação são específicos para o  \textit{assembler} \textbf{yasm}. Outros montadores podem
ser um pouco diferente. Um programa completo na linguagem assembly é apresentado para demonstrar a formatação apropriada do programa.

Um arquivo de origem de montagem formatado corretamente consiste em várias partes principais:
\begin{itemize}
	\item Seção de dados, na qual os dados inicializados são declarados e definidos.
	\item Seção BSS, na qual dados não inicializados são declarados.
	\item Seção de texto, na qual o código é colocado.
\end{itemize}

As seções a seguir resumem os requisitos básicos de formatação. Somente o básico da formatação e sintaxe do assembler serão apresentadas. Para informações adicionais, consulte o
manual de referência \textbf{yasm} (como observado no Capítulo \ref{cap1}, Introdução).

\section{Comentários}
O ponto-e-vírgula (\textbf{;}) é usado para anotar os comentários do programa. Comentários (usando o \textbf{;}) podem ser
colocados em qualquer lugar, inclusive após uma instrução. Qualquer caractere após o \textbf{;} é ignorado pelo montador. Isso pode ser usado para explicar as etapas executadas no código ou para comentar
seções de código.

\section{Valores numéricos}
Os valores numéricos podem ser especificados em decimal, hexadecimal ou octal. 

Ao especificar valores hexadecimais ou de base 16, eles são precedidos por um \textbf{0x}. Por exemplo, para especificar \textbf{127} como hexadecimal, usamos \textbf{0x7f}.

Ao especificar valores octais, ou base 8, eles são seguidos por um \textbf{q}. Por exemplo, para especificar \textbf{511} como octal, seria \textbf{777q}.

O radical padrão (base) é decimal, portanto, nenhuma notação especial é necessária para números decimais (base 10).

\section{Definindo constantes}
Constantes são definidas com \textbf{equ}. O formato geral é:

\begin{verbatim}
<nome> equ <valor>
\end{verbatim}

O valor de uma constante não pode ser alterado durante a execução do programa.

As constantes são substituídas pelos valores definidos durante o processo de montagem. Sendo assim, uma constante não recebe um local de memória. Isso torna a constante mais flexível, pois não é atribuído um tipo/tamanho específico (byte, palavra, palavra dupla, etc.). Os valores estão sujeitos às limitações de alcance do uso pretendido. Por exemplo, a seguinte constante,
\begin{verbatim}
TAMANHO equ 10000
\end{verbatim}
pode ser usada como uma word ou word dupla, mas não como um byte.

\section{Seção de dados}
Os dados inicializados devem ser declarados na seção ``\textit{section .data}''. Deve haver um espaço após a palavra `\textit{section}'. Todas as variáveis e constantes são colocadas nesta seção. Os nomes de variáveis devem começar com uma letra, seguida de letras ou números, incluindo alguns caracteres especiais (como o sublinha, ``\_''). Definições de variáveis devem incluir o nome, o tipo de dados e o valor inicial da variável.

O formato geral é:
\begin{verbatim}
<nomeVariável> <tipoDeDado> <valorInicial>
\end{verbatim}

Consulte as seções a seguir para uma série de exemplos usando vários tipos de dados.

Os tipos de dados suportados estão na Tabela \ref{tab:tiposDeDados}.

% Please add the following required packages to your document preamble:
% \usepackage{booktabs}
% \usepackage[table,xcdraw]{xcolor}
% If you use beamer only pass "xcolor=table" option, i.e. \documentclass[xcolor=table]{beamer}
\begin{table}[]
	\centering
	\begin{tabular}{|c|l|}
		\hline
		\rowcolor[HTML]{C0C0C0} 
		{\color[HTML]{000000} } Declaração & {\color[HTML]{000000} } \\ \hline
		\textbf{db}& variáveis de 8 bits \\ \hline
		\textbf{dw} & variáveis de 16 bits\\ \hline
		\textbf{dd} & variáveis de 32 bits\\ \hline
		\textbf{dq} & variáveis de 64 bits\\ \hline
		\textbf{ddq} & variáveis de 128 bits $\rightarrow$ inteira\\ \hline
		\textbf{dt} & variáveis de 128 bits $\rightarrow$ ponto flutuante\\ \hline
	\end{tabular}
\caption{Tipos de dados.}
\label{tab:tiposDeDados}
\end{table}

Estas são as diretivas principais do assembler para declarações de dados inicializados. Outras diretivas são referenciadas em diferentes seções.

Vetores inicializadas são definidos com valores separados por vírgula.

Alguns exemplos simples incluem:

\begin{verbatim}
bVar    db  10              ; variável byte
cVar    db  "H"             ; caracter único
strng   db  "Hello World"   ; string
wVar    dw 	5000            ; variável word
dVar    dd 	50000           ; variável de 32 bits
arr     dd 	100, 200, 300   ; array com 3 elementos
flt1    dd 	3.14159         ; ponto flutuante de 32 bits
qVar    dq 	1000000000      ; variável de 64 bits
\end{verbatim}

O valor especificado deve poder caber no tipo de dados especificado. Por exemplo, se o valor de uma variável de tamanho de byte é definido como 500, geraria um erro do assembler.

\section{Seção BSS}
Dados não inicializados são declarados na seção ``section .bss'' (Block Started by Symbol). Deve haver um espaço depois da palavra 'section'. Todas as 
variáveis não inicializadas são declaradas nesta seção. Os nome de variável começam com uma letra seguida de letras ou números, incluindo alguns caracteres especiais (como o sublinha, ``\_''). As definições de variáveis devem incluir o tipo de dados e a contagem.

O formato geral é:
\begin{verbatim}
<nomeDeVariável> <resTipo> <contador>
\end{verbatim}

Consulte as seções a seguir para uma série de exemplos usando vários tipos de dados.

Os tipos de dados suportados estão na Tabela \ref{tabDadosNInicializados}:

{\centering
\begin{table}[ht]
	\centering
	\begin{tabular}{|c|l|}
		\hline
		\rowcolor[HTML]{C0C0C0} 
		{\color[HTML]{000000} } Declaração & {\color[HTML]{000000} } \\ \hline
		\textbf{resb}& variáveis de 8 bits \\ \hline
		\textbf{resw} & variáveis de 16 bits\\ \hline
		\textbf{resd} & variáveis de 32 bits\\ \hline
		\textbf{resq} & variáveis de 64 bits\\ \hline
		\textbf{resdq} & variáveis de 128 bits \\ \hline
	\end{tabular}
    \caption{Declarações de variáveis não inicializadas}
    \label{tabDadosNInicializados}
\end{table}
} 

Essas são as diretivas do assembler principal para declarações de dados não inicializados. Outras diretivas são referenciadas em diferentes seções.

Alguns exemplos simples incluem:

\begin{verbatim}
bArr  resb 10  ; vetor com 10 elementos com tamanho byte
wArr  resw 50  ; vetor com 50 elementos com tamanho word
dArr  resd 100 ; vetor com 100 elementos com tamanho double
qArr  resq 200 ; vetor com 10 elementos com tamanho quad
\end{verbatim}

A matriz alocada não é inicializada para nenhum valor específico.

\section{Seção de texto}
O código é colocado na seção ``section .text''. Deve haver um espaço após a palavra 'section'. As instruções são especificadas uma por linha e cada uma deve ser uma instrução válida com os operandos necessários apropriados.

A seção de texto incluirá alguns cabeçalhos ou etiquetas que definem o ponto de entrada inicial do programa. Por exemplo, assumindo um programa básico usando o \textit{linker} padrão do sistema, as seguintes declarações devem ser incluídas.
\begin{verbatim}
global _start
_start:
\end{verbatim}

Nenhuma etiqueta ou diretiva especial é necessária para encerrar o programa. No entanto, um serviço do sistema deve ser usado para informar ao sistema operacional que o programa deve ser encerrado e os recursos, como memória, recuperados e reutilizados. Consulte o
exemplo de programa na seção a seguir.

\section{Programa de exemplo}
Um programa em linguagem assembly muito simples é apresentado para demonstrar a formatação do programa.

\begin{verbatim}
; Exemplo simples que demonstra o formato e o layout básicos do programa.

; Ed Jorgensen
; July 18, 2014

; ************************************************************
; Algumas declarações básicas de dados

section .data

; -----
; Define constantes

EXIT_SUCCESS      equ 0  ; operação bem sucedida
SYS_exit          equ 60 ; chamada de código para terminar

; -----
; Declarações de variáveis Byte (8-bits)

bVar1             db 17 
bVar2             db 9
bResult           db 0

; -----
; Declarações de variáveis word (16-bits)
wVar1             dw 17000
wVar2             dw 9000
wResult           dw 0

; -----
; Declarações de variáveis word dupla (32-bits)
dVar1             dd 17000000
dVar2             dd 9000000
dResult           dd 0

; -----
; Declarações de variáveis quad-word dupla (64-bits)
qVar1             dq 170000000
qVar2             dq 90000000
qResult           dq 0

; ************************************************************
; Seção de código
section .text
global _start
_start:

; ----------
; Executa uma série de operações de adição muito básicas
; para demonstrar o formato básico do programa.
; ----------
; Exemplo de byte
; bResult = bVar1 + bVar2
    mov    al, byte [bVar1]
    add    al, byte [bVar2]
    mov    byte [bResult], al

; ----------
; Exemplo word
; wResult = wVar1 + wVar2
    mov   ax, word [wVar1]
    add   ax, word [wVar2]
    mov   word [wResult], ax

; ----------
;  double-word
; dResult = dVar1 + dVar2
    mov   eax, dword [dVar1]
    add   eax, dword [dVar2]
    mov   dword [dResult], eax

; ----------
; Exemplo quad-word
; qResult = qVar1 + qVar2
    mov   rax, qword [qVar1]
    add   rax, qword [qVar2]
    mov   qword [qResult], rax

; ************************************************************
; Feito, termina o programa.
last:
    mov rax, SYS_exit
    mov rdi, EXIT_SUCCESS
    syscall
\end{verbatim}

Este exemplo de programa será referenciado e melhor explicado nos capítulos seguintes.

\section{Exercícios}
Abaixo estão algumas questões baseadas neste capítulo.

\subsection{Perguntas}
Abaixo estão algumas perguntas do questionário.
\begin{enumerate}
	\item Qual é o nome do montador sendo usado neste capítulo?
	\item Como os comentários são marcados em um programa em linguagem assembly?
	\item Qual é o nome da seção em que os dados inicializados foram declarados?
	\item Qual é o nome da seção em que os dados não inicializados foram declarados?
	\item Qual é o nome da seção onde o código é colocado?
	\item Qual é a declaração de dados para cada uma das seguintes variáveis com os valores fornecidos:
	\begin{enumerate}
		\item  variável de tamanho byte \textbf{\textit{bNum}} definida como $10_{10}$
		\item variável de tamanho de palavra \textbf{\textit{wNum}} definida como $10.291_{10}$
		\item variável \textbf{\textit{dwNum}} de tamanho de palavra dupla definida como $2.126.010_{10}$
		\item variável de tamanho quádruplo \textbf{\textit{qwNum} } definida como $10.000.000.000_{10}$
	\end{enumerate}
    \item Qual é a declaração de dados não inicializados para cada um dos seguintes casos:    
    \begin{enumerate}
    	\item matriz com tamanho de bytes chamada \textbf{\textit{bArr}} com 100 elementos
    	\item array com tamanho de palavra chamado \textbf{\textit{wArr} }com 3000 elementos
    	\item array com tamanho double-word, chamado \textbf{\textit{dwArr}}, com 200 elementos
    	\item array com tamanho de quad-word chamado \textbf{\textit{qArr}} com 5000 elementos
    \end{enumerate}
    \item Quais são as declarações necessárias para indicar o início de um programa (na seção texto)?

\end{enumerate}



\chapter{Ferramentas de desenvolvimento}
\label{Ferramentasdedesenvolvimento}

\epigraph{``Existem duas maneiras de escrever programas sem erros; somente a terceira funciona.''}


Em geral, o conjunto de ferramentas de programação usadas para criar um programa é chamado de ferramentas de desenvolvimento. Para os fins deste texto, a cadeia de ferramentas consiste no seguinte:
\begin{itemize}
	\item Assembler (montador)
	\item Linker (ligador)
	\item Loader (carregador)
	\item Debugger (depurador)
\end{itemize}

Embora existam muitas opções para a cadeia de ferramentas, este texto usa um conjunto bastante padrão de ferramentas de código aberto que funcionam bem juntas e oferecem suporte total ao ambiente x86 de 64 bits.

Cada uma dessas ferramentas de programação é explicada nas seções a seguir.
\section{Visão geral de montagem/link/carga}
Em termos gerais, o processo de montar, \textit{linkar} e carregar é como os arquivos de origem escritos pelo programador são convertidos em um programa executável.

O arquivo de origem legível por humanos é convertido em um arquivo-objeto pelo assembler. Na forma mais básica, o arquivo-objeto é convertido em um arquivo executável pelo linker. O carregador carregará o arquivo executável na memória.

Uma visão geral do processo é fornecida no diagrama a seguir.

\tikzstyle{decision} = [diamond, draw, fill=blue!20, 
text width=4.5em, text badly centered, node distance=3cm, inner sep=0pt]
\tikzstyle{block} = [rectangle, draw, fill=blue!20, 
text width=5em, text centered, rounded corners, minimum height=4em]
\tikzstyle{bin} = [rectangle, draw, fill=blue!40, 
text width=5em, text centered, minimum height=4em]
\tikzstyle{line} = [draw, -latex']
\tikzstyle{cloud} = [draw, ellipse,fill=red!20, node distance=3cm,
minimum height=2em]

\tikzstyle{doc}=[%
draw,
thick,
align=center,
color=black,
shape=document,
minimum width=20mm,
shape=document,
inner sep=2ex,
]

\begin{tikzpicture}[node distance = 2.5cm, auto]
% Place nodes
\node [block] (init) {Assembler};
\node[doc] at ([shift=(0:-3cm)]init) (assembly) {Fonte em\\ linguagem\\ assembly};
\node [bin, below of=init] (objeto) {Arquivo objeto};
\node [block, below of=objeto] (link) {Link};
\node [doc, left of=objeto, node distance=3cm] (listagem) {Arquivo\\listagem};
\node [bin, node distance=3cm] at ([shift=(5:2.5cm)]objeto) (outro) {Outro arquivo objeto (se houver)};
\node [bin, node distance=3cm] at ([shift=(-4:3cm)]outro) (biblio) {Rotinas da biblioteca (se houver)};
\node [bin, below of=link] (exe) {Executável};
\node [bin, left of=exe] (shared) {Arquivos de objetos compartilhados};
\node [block, right of=exe, node distance=3cm] (load) {Load};
\node [block, below of=exe, node distance=3cm] (ram) {RAM};
% Draw edges
\path [line] (assembly) -- (init);
\path [line] (init) -- (objeto);
\path [line] (init) -- (listagem);
\path [line] (objeto) -- (link);
\path [line, dashed] (outro) -- (link);
\path [line, dashed] (biblio) |- (link);

\path [line] (link) -- (exe);
\path [line] (exe) -- (load);
\path [line] (load) -- (ram);
\path [line, dashed] (shared) -- (ram);


\end{tikzpicture}

As etapas de montagem, ligação e carregamento são descritas com mais detalhes nas seções a seguir.

\section{Assembler (montador)}
O assembler é um programa que lê um arquivo de entrada em linguagem \textit{assembly} e converte o código em um arquivo binário de linguagem de máquina. O arquivo de entrada é um arquivo fonte em linguagem \textit{assembly} que contém instruções em linguagem \textit{assembly} em formato legível por humanos. A saída da linguagem de máquina é conhecida como um arquivo de objeto. Como parte desse processo, os comentários são removidos e os nomes das variáveis e rótulos são convertidos em endereços apropriados (conforme exigido pela CPU durante a execução).

O assembler usado neste texto é o \textbf{yasm} 25. Links para o site e documentação do \textbf{yasm} podem ser encontrados no Capítulo \ref{cap1}, Introdução.

\subsection{Comandos de montagem}
O comando do \textbf{yasm} apropriado para ler o arquivo de origem da linguagem \textit{assembly}, como o exemplo do capítulo anterior, é o seguinte:
\begin{center}
	\textbf{yasm -g dwarf2 -f elf64 exemplo.asm -l exemplo.lst}
\end{center}

Observe que o \textit{-l} é uma letra \textbf{L} minúscula com traço na frente (que pode ser facilmente confundida com o número $ 1$).

A opção \textbf{-g dwarf2} é usada para informar ao montador para incluir informações de depuração no arquivo de objeto final. Isso aumenta o tamanho do arquivo de objeto, mas é necessário para permitir uma depuração eficaz. O \textbf{-f elf64} informa o montador para criar o arquivo objeto no formato \textbf{ELF64}, que é apropriado para sistemas baseados em \textit{Linux de 64 bits}. O \textbf{exemplo.asm} é o nome do arquivo de origem da linguagem \textit{assembly} para entrada. O \textbf{-l exemplo.lst} (traço minúsculo, letra L) informa ao montador para criar um arquivo de lista denominado \textit{exemplo.lst}.

Se ocorrer um erro durante o processo de montagem, ele deve ser resolvido antes de prosseguir para a etapa de link.

\subsection{Arquivo de listagem}
Além disso, o montador é opcionalmente capaz de criar um arquivo de listagem. O arquivo de listagem mostra o número da linha, o endereço relativo, a versão em linguagem de máquina da instrução (incluindo referências de variáveis) e a linha de origem original. O arquivo de listagem pode ser útil durante a depuração.

Por exemplo, um fragmento da seção de dados do arquivo de lista, do programa de exemplo no capítulo anterior, é o seguinte:

\begin{verbatim}
36 00000009 40660301	dVar1	  dd	17000000
37 0000000D 40548900	dVar2	  dd	9000000
38 00000011 00000000	dResult	dd	0
\end{verbatim}

Na primeira linha, \textbf{36} é o número da linha. O próximo número, \textbf{0x00000009}, é o endereço relativo na área de dados de onde a variável será armazenada. Como \textbf{dVar1} é uma palavra dupla, que requer quatro bytes, o endereço da próxima variável é \textbf{0x0000000D}. A variável \textbf{dVar1} usa \textbf{4} bytes como endereços \textbf{0x00000009}, \textbf{0x0000000A}, \textbf{0x0000000B} e \textbf{0x0000000C}. O resto da linha é a declaração de dados conforme digitada no arquivo-fonte original da linguagem \textit{assembly}.

O \textbf{0x40660301} é o valor, em hexadecimal, conforme colocado na memória. O \textbf{$ 17000000_{10} $}  é \textbf{0x01036640}. Lembrando que a arquitetura é \textit{little-endian}, o byte menos significativo (\textbf{0x40}) é colocado no endereço de memória mais baixo. Como tal, o \textbf{0x40} é colocado no endereço relativo \textbf{0x00000009}, o próximo byte, \textbf{0x66}, é colocado no endereço \textbf{0x0000000A} e assim por diante. Isso pode ser confuso, pois à primeira vista o número pode aparecer invertido ou truncado (dependendo de como é visualizado).

Para ajudar a visualizar, a imagem da memória seria a a da Figura \ref{fig:litllelayout}:

\tikzstyle{memoria} = [rectangle, draw, fill=blue!20, 
text width=4em, text centered]

\begin{figure}[h]
\begin{center}
		\begin{tikzpicture}
\node [text width=5em,xshift=0.5cm,] (nome) {\textbf{Nome da variável}};
\node [text width=5em, xshift=3cm, yshift=-0.2cm] (valor)  {\textbf{Valor}};
\node [text width=5em, xshift=4.5cm, yshift=-0.2cm] (endereco) {\textbf{Endereço}};
\node [memoria, xshift=2.5cm, yshift=-0.8cm ]  {00};
\node [xshift=4.5cm, yshift=-0.8cm ]  {0x00000010};
\node [memoria, xshift=2.5cm, yshift=-1.3cm ]  {89};
\node [xshift=4.5cm, yshift=-1.3cm ]  {0x0000000F};
\node [memoria, xshift=2.5cm, yshift=-1.8cm ]  {54};
\node [xshift=4.5cm, yshift=-1.8cm ]  {0x0000000E};
\node [memoria, xshift=2.5cm, yshift=-2.3cm ]  {40};
\node [xshift=0.5cm, yshift=-2.3cm ]  {dVar2 $\Rightarrow$};
\node [xshift=4.5cm, yshift=-2.3cm ]  {0x0000000D};
\node [memoria, xshift=2.5cm, yshift=-2.8cm ]  {01};
\node [xshift=4.5cm, yshift=-2.8cm ]  {0x0000000C};
\node [memoria, xshift=2.5cm, yshift=-3.3cm ]  {03};
\node [xshift=4.5cm, yshift=-3.3cm ]  {0x0000000B};
\node [memoria, xshift=2.5cm, yshift=-3.8cm ]  {66};
\node [xshift=4.5cm, yshift=-3.8cm ]  {0x0000000A};
\node [memoria, xshift=2.5cm, yshift=-4.3cm ]  {40};
\node [xshift=4.5cm, yshift=-4.3cm ]  {0x00000009};
\node [xshift=0.5cm, yshift=-4.3cm ]  {dVar1 $\Rightarrow$};

\end{tikzpicture}
\end{center}
	\caption{Little-Endian, Layout de dados de múltiplas variáveis}
	\label{fig:litllelayout}
\end{figure}

Por exemplo, um fragmento da seção de texto do arquivo de lista, extraído do programa de exemplo no capítulo anterior, é o seguinte:

\begin{verbatim}
95                         last:
96 0000005A 48C7C03C000000	  mov		rax, SYS_exit
97 00000061 48C7C300000000	  mov		rdi, EXIT_SUCCESS	
98 00000068 0F05			            syscall
\end{verbatim}

Novamente, os números à esquerda são os números das linhas. O próximo número, \textbf{0x0000005A}, é o endereço relativo de onde a linha de código será colocada.

O próximo número, \textbf{0x48C7C03C000000}, é a versão em linguagem de máquina da instrução, em hexadecimal, que a CPU lê e entende. O resto da linha é a instrução de origem da linguagem \textit{assembly} original.

O rótulo, \textbf{last:}, não tem uma instrução em linguagem de máquina, pois o rótulo é usado para fazer referência a um endereço específico e não é uma instrução executável.

\subsection{Montador de dois passos}
O montador lerá o arquivo fonte e converterá cada instrução da linguagem \textit{assembly}, digitada pelo programador, em um conjunto de uns e zeros que a CPU sabe que é essa instrução. Os 1s e 0s são chamados de \textbf{linguagem de máquina}. Há uma correspondência de um para um entre as instruções em linguagem \textit{assembly} e a \textbf{linguagem de máquina binária}. Esse relacionamento significa que a linguagem de máquina, na forma de um arquivo executável, pode ser convertida novamente em linguagem \textit{assembly} legível por humanos. Claro, os comentários, nomes de variáveis e nomes de rótulos estão faltando, então o código resultante pode ser muito difícil de ler.

Conforme o montador lê cada linha da linguagem \textit{assembly}, ele gera um código de máquina para essa instrução. Isso funcionará bem para instruções que não realizam saltos. No entanto, para instruções que podem alterar o fluxo de controle (por exemplo, instruções \textbf{IF}, saltos incondicionais), o montador não é capaz de converter a instrução. Por exemplo, dado o seguinte fragmento de código:

\begin{verbatim}
       mov rax, 0
       jmp skipRest
       ...
       ...
    skipRest:
\end{verbatim}

Isso é conhecido como uma \textbf{referência antecipada}. Se o montador lê o arquivo de montagem uma linha por vez, ele não leu a linha onde \textbf{skipRest} está definido. Na verdade, ele nem sabe ao certo se \textbf{skipRest} está definido.

Esta situação pode ser resolvida lendo o arquivo de origem do \textit{assembly} duas vezes. Todo o processo é conhecido como \textbf{assembler} de duas passagens. As etapas necessárias para cada passagem são detalhadas nas seções a seguir.

\subsubsection{Primeiro passo}
As etapas realizadas na primeira passagem variam com base no projeto do montador específico. No entanto, algumas das operações básicas realizadas na primeira passagem incluem o seguinte:

\begin{itemize}
	\item Criar tabela de símbolos
	\item Expandir macros
	\item Avaliar expressões constantes
\end{itemize}

Uma macro é um elemento de programa que é expandido em um conjunto de instruções predefinidas pelo programador. Para obter mais informações, consulte o Capítulo \ref{cap:Macros}, Macros.

Uma expressão constante é uma expressão composta inteiramente de constantes. Como a expressão é apenas constante, ela pode ser totalmente avaliada no momento da montagem. Por exemplo, supondo que a constante \textbf{BUFF} seja definida, a seguinte instrução contém uma expressão constante;
\begin{verbatim}
mov rax, BUFF+5
\end{verbatim}

Este tipo de expressão constante é comumente usado em programas grandes ou complexos.

Os endereços são atribuídos a todas as instruções do programa. A tabela de símbolos é uma lista ou tabela de todos os símbolos de programa, nomes de variáveis e rótulos de programa e seus respectivos endereços no programa.

Conforme apropriado, algumas diretivas do montador são processadas na primeira passagem.

\subsubsection{Segundo passo}
As etapas executadas na segunda passagem variam com base no projeto do montador específico. No entanto, algumas das operações básicas realizadas na segunda passagem incluem o seguinte:

\begin{itemize}
	\item Geração final de código
	\item Criação de arquivo de lista (se solicitado)
	\item Criação de arquivo-objeto
\end{itemize}

O termo \textit{geração de código} refere-se à conversão da instrução de linguagem de montagem fornecida pelo programador na instrução de linguagem de máquina executável da CPU. Devido à correspondência um para um, isso pode ser feito para instruções que não usam símbolos na primeira ou na segunda passagem.

Deve-se observar que, com base no design do montador, grande parte da geração do código pode ser feita na primeira passagem ou toda feita na segunda passagem. De qualquer forma, a geração final é executada na segunda passagem. Isso exigirá o uso da tabela de símbolos para verificar os símbolos do programa e obter os endereços apropriados da tabela.

O arquivo de lista, embora opcional, pode ser útil para depuração. Se solicitado, ele seria gerado na segunda passagem.

Se não houver erros, o arquivo de objeto final será criado na segunda passagem.

\subsection{Diretivas do Assembler}
As diretivas do \textit{assembler} são instruções para o montador que o direcionam a fazer algo. Isso pode ser formatação ou layout. Essas diretivas não são traduzidas em instruções para a CPU.

\section{Linker}
O \textbf{linker}, às vezes referido como editor de ligação, combinará um ou mais arquivos-objeto em um único arquivo executável. Além disso, todas as rotinas das bibliotecas do usuário ou do sistema são incluídas conforme necessário. O \textbf{linker GNU gold, ld}, é usado. O comando do vinculador apropriado para o programa de exemplo do capítulo anterior é o seguinte:
\begin{center}
	\textbf{ld -g -o exemplo exemplo.o}
\end{center}

Observe que o \textbf{-o} é um traço minúsculo, letra \textbf{O}, que pode ser confundido com o número \textbf{0}.

A opção \textbf{-g} é usada para informar ao vinculador para incluir informações de depuração no arquivo executável final. Isso aumenta o tamanho do arquivo executável, mas é necessário para permitir uma depuração eficaz. O  \textbf{-o exemplo} especifica a criação do arquivo executável denominado \textbf{exemplo} (sem extensão). Se a opção \textbf{-o <fileName>} for omitida, o arquivo de saída será denominado \textbf{a.out} (por padrão). O \textbf{exemplo.o} é o nome do arquivo de objeto de entrada lido pelo vinculador. Deve-se observar que o arquivo executável pode ter qualquer nome e não precisa ter o mesmo nome de qualquer um dos arquivos de objeto de entrada.

\subsection{Vinculando vários arquivos}
Na programação, grandes problemas são normalmente resolvidos dividindo-os em problemas menores. Os problemas menores podem ser resolvidos individualmente, possivelmente por diferentes programadores.

Arquivos de objetos de entrada adicionais, se houver, seriam listados, em ordem, separados por um espaço. Por exemplo, se houver dois arquivos de objeto, \textbf{main.o} e \textbf{funcs.o}, o comando de link para criar um exemplo de nome de arquivo executável, com informações de depuração incluídas, seria o seguinte:
\begin{center}
	\textbf{ld -g -o exemplo main.o funcs.o}
\end{center}

Isso normalmente seria necessário para programas maiores ou mais complexos.

Ao usar funções localizadas em um arquivo de origem externa diferente, qualquer função ou funções que não estejam no arquivo de origem atual devem ser declaradas como externas. Variáveis, como variáveis globais, em outros arquivos de origem também podem ser acessadas usando a instrução \textbf{extern}, no entanto, os dados são normalmente transferidos como argumentos da chamada de função.

\subsection{Processo de Ligação}
A ligação é o processo fundamental de combinar as soluções menores em uma única unidade executável. Se qualquer usuário ou rotinas de biblioteca do sistema forem usadas, o \textit{linker} incluirá as rotinas apropriadas. Os arquivos de objeto e rotinas de biblioteca são combinados em um único módulo executável. O código da linguagem de máquina é copiado de cada arquivo de objeto em um único executável.

Como parte da combinação dos arquivos de objeto, o \textit{linker} deve ajustar os endereços realocáveis conforme necessário. Supondo que haja dois arquivos de origem, o principal e um arquivo de origem secundário contendo algumas funções, ambos reunidos em arquivos de objeto \textbf{main.o} e \textbf{funcs.o}. Quando cada arquivo é montado, as chamadas para rotinas fora do arquivo que está sendo montado são declaradas com a diretiva de assembler \textbf{extern}. O código não está disponível para uma referência externa e tais referências são marcadas como externas no arquivo de objeto. O arquivo de lista mostrará um \textbf{R} para esses endereços relocáveis. O \textit{linker} deve satisfazer as referências externas. Além disso, a localização final das referências externas deve ser colocada no código.

Por exemplo, se o arquivo de objeto \textbf{main.o} chamar uma função no arquivo \textbf{funcs.o}, o \textit{linker} deverá atualizar a chamada com o endereço apropriado, conforme mostrado na Figura \ref{fig:linker}.

\begin{figure}[h]
\begin{center}
	\begin{tikzpicture}[node distance = 4cm, auto]
% Place nodes
\node[doc, fill=blue!20] (init) {main.o\\$\ldots$\\\\\\\ldots\\};
\node [doc, fill=blue!20, below of=init] (objeto) {funcs.o\\$\ldots$\\\\0x0100:\\$\ldots$};
\node [doc,fill=blue!20, right of=init] (exe) {executável\\\\call 0x0400\vspace{1cm}};
\node [bin, below of=exe, thick, text width=2.31cm,  fill=blue!20,node distance=3cm] (objeto) {\vspace{0.8cm}\\$\ldots$\\0x400:\\$\ldots$\vspace{1cm}};
\node at (0,-0.1) (call) {call fnc1};
\node at (-0.2,-4) (fnc) {fnc1:};
\draw[thick, ->] (1,-0.1) -- (2.7,-0.1);
\draw[thick, ->] (1,-4) -- (2.7,-3.1);

\end{tikzpicture}
\end{center}
	\caption{Ligando vários arquivos}
	\label{fig:linker}
\end{figure}

Aqui, a função \textbf{fnc1} é externa ao arquivo de objeto \textbf{main.o} e é marcada com um \textbf{R}. A função real \textbf{fnc1} está no arquivo \textbf{funcs.o}, que inicia seu endereçamento relativo de \textbf{0x0} (na seção de texto), pois ela não sabe nada sobre o código principal. Quando os arquivos de objeto são combinados, o endereço relativo original de \textbf{fnc1} (mostrado como \textbf{0x0100:}) é alterado para seu endereço final no arquivo executável (mostrado como \textbf{0x0400:}). O \textit{linker} deve inserir este endereço final na instrução de chamada no principal (mostrado como chamada \textbf{0x0400:}) para concluir o processo de vinculação e garantir que a chamada de função funcione corretamente.

Isso ocorrerá com todos os endereços realocáveis para código e dados.

\subsection{Linking Dinâmico}
O sistema operacional Linux suporta links dinâmicos, o que permite adiar a resolução de alguns símbolos até que um programa esteja sendo executado. As instruções reais não são colocadas no arquivo executável e, em vez disso, se necessário, são resolvidas e acessadas em tempo de execução.

Embora mais complexa, essa abordagem oferece duas vantagens:\begin{itemize}
	\item Bibliotecas frequentemente utilizadas (por exemplo, as bibliotecas do sistema padrão) podem ser armazenadas em apenas um local, não duplicadas em cada binário único.
	\item Se um \textit{bug} em uma função de biblioteca for corrigido, todos os programas que o usam dinamicamente se beneficiarão da correção (na próxima execução). Caso contrário, os programas que utilizam esta função por vinculação estática teriam que ser vinculados novamente antes de aplicar a correção.
\end{itemize}

Também existem desvantagens:
\begin{itemize}
	\item Uma biblioteca atualizada incompatível interromperá os executáveis que dependiam do comportamento da versão anterior da biblioteca.
	\item Um programa, junto com as bibliotecas que usa, pode ser certificado (por exemplo, quanto à exatidão, requisitos de documentação ou desempenho) como um pacote, mas não se os componentes puderem ser substituídos.
\end{itemize}

No Linux/Unix, os arquivos de objeto ligados dinamicamente geralmente têm a extensão \textbf{.so} (\textit{shared object}, objeto compartilhado). No \textit{Windows}, eles têm uma extensão \textbf{.dll} (\textit{dynamically linked library}, biblioteca vinculada dinamicamente). Mais detalhes sobre links dinâmicos estão fora do escopo deste texto.

\section{Script para montagem/ligação}
Ao programar, geralmente é necessário digitar os comandos de montagem e link muitas vezes com vários programas diferentes. Em vez de digitar os comandos de montagem (\textbf{yasm}) e link (\textbf{ld}) a cada vez, é possível colocá-los em um arquivo, denominado arquivo de \textit{script}. Em seguida, pode-se executar o arquivo de \textit{script} que executará apenas os comandos que foram inseridos no arquivo. Embora não seja obrigatório, o uso de um arquivo de \textit{script} pode economizar tempo e tornar as coisas mais fáceis ao trabalhar em um programa.

Um exemplo simples de script no shell bash para montagem e link é o seguinte:
\begin{verbatim}
#!/bin/bash
# Simple assemble/link script.
if [ -z $1 ]; then
   echo "Usage: ./asm64 <asmMainFile> (no extension)"
   exit
fi
# Verifique se nenhuma extensão foi inserida
if [ ! -e "$1.asm" ]; then
   echo "Error, $1.asm not found."
   echo "Note, do not enter file extensions."
   exit
fi
# Compile, assemble e link.

yasm -Worphan-labels -g dwarf2 -f elf64 $1.asm -l $1.lst
ld -g -o $1 $1.o
\end{verbatim}

O \textit{script} acima deve ser colocado em um arquivo. Para este exemplo, o arquivo será denominado \textbf{asm64} e colocado no diretório de trabalho atual (onde os arquivos de origem estão localizados).

Uma vez criado, será necessário o privilégio de execução para o arquivo de \textit{script}, como segue:
\begin{center}
	\begin{verbatim}
chmod +x asm64
\end{verbatim}
\end{center}

Isso só precisará ser feito uma vez para cada arquivo de \textit{script}.

O arquivo de \textit{script} lerá o nome do arquivo de origem na linha de comando. Por exemplo, para usar o arquivo de \textit{script} para montar o exemplo do capítulo anterior (denominado \textbf{exemplo.asm}), digite o seguinte:
\begin{center}
	\begin{verbatim}
	./asm64 exemplo
	\end{verbatim}
\end{center}

A extensão ``\textbf{.asm}'' no arquivo \textit{exemplo.asm} não deve ser incluída (uma vez que é adicionada no \textit{script}). O arquivo de \textit{script} montará e vinculará o arquivo de origem, criando o arquivo de lista, o arquivo de objeto e o arquivo executável.

O uso deste ou de qualquer arquivo de \textit{script} é opcional. O nome do arquivo de \textit{script} pode ser alterado conforme desejado.

\section{Loader}
O carregador (\textit{loader}) é uma parte do sistema operacional que irá carregar o programa do armazenamento secundário para o armazenamento primário (isto é, memória principal). Em termos gerais, o carregador tentará encontrar e, se encontrado, ler um arquivo executável formatado corretamente, criar um novo processo, carregar o código na memória e marcar o programa como pronto para execução. O \textit{scheduler} do sistema operacional tomará as decisões sobre qual processo é executado e quando o processo é executado.

O \textit{loader} é invocado implicitamente digitando o nome do programa. Por exemplo, no programa de exemplo anterior, denominado \textbf{exemplo}, o comando \textit{Linux} seria:
\begin{center}
	\begin{verbatim}
	./exemplo
	\end{verbatim}
\end{center}
que executará o arquivo denominado exemplo criado pelas etapas anteriores (montagem e vinculação). Como o programa de exemplo não executa nenhuma saída, nada será exibido no console. Como tal, um depurador pode ser usado para verificar os resultados.

\section{Debugger}
O depurador (\textit{debugger}) é usado para controlar a execução de um programa. Isso permite que as atividades de teste e depuração sejam realizadas.

No exemplo anterior, o programa computou uma série de cálculos, mas não produziu nenhum dos resultados. O depurador pode ser usado para verificar os resultados. O arquivo executável é criado com o comando de montagem e link descrito anteriormente e deve incluir a opção \textbf{-g}.

O depurador usado é o depurador \textbf{GNU DDD} que fornece uma interface visual para o depurador de linha de comando GNU, \textbf{gdb}. O site e a documentação do \textbf{DDD} são mencionados na seção de referências do Capítulo \ref{cap1}, Introdução.

Devido à complexidade e importância do depurador, um capítulo separado para a depuração é fornecido.

\section{Exercícios}
Abaixo estão algumas perguntas do questionário.
\begin{enumerate}
	\item Qual é a relação entre linguagem assembly e linguagem de máquina?
	\item Quais ações são realizadas na primeira passagem do montador?
	\item Quais ações são executadas na segunda passagem do montador?
	\item Quais ações são realizadas pelo linker?
	\item Quais ações são realizadas pelo loader?
	\item Forneça um exemplo de uma expressão constante.
	\item Desenhe um diagrama de todo o processo de montagem, link e carregamento.
	\item Quando um arquivo de objeto compartilhado é vinculado a um programa?
	\item O que está contido na tabela de símbolos (duas coisas)?
\end{enumerate}



\chapter{Debugger DDD}

\epigraph{``Meu software nunca possui bugs. Ele só
	desenvolve recursos aleatórios.''}

Um depurador permite ao usuário controlar a execução de um programa, examinar variáveis, outra memória (ou seja, espaço na pilha) e exibir a saída do programa (se houver). O GNU Data Display Debugger de código aberto (DDD\footnote{Para mais informações, consulte: https://pt.wikipedia.org/wiki/Data\_Display\_Debugger}) é um front-end visual para o GNU Debugger (GDB\footnote{Para mais informações, consulte: https://pt.wikipedia.org/wiki/GNU\_Debugger}) e está amplamente disponível. Outros depuradores podem ser facilmente usados, se desejado.

Somente os comandos básicos do depurador são abordados neste capítulo. O depurador DDD possui muitos outros recursos e opções não abordados aqui. À medida que você ganha experiência, vale a pena revisar a documentação do DDD, mencionada no Capítulo \ref{cap1}, para aprender mais sobre os recursos adicionais, a fim de ajudar a melhorar a eficiência geral da depuração.

A funcionalidade do DDD pode ser estendida usando vários plug-ins. Os plug-ins não são necessários e não serão abordados neste capítulo.

Este capítulo aborda o uso do depurador GNU DDD como uma ferramenta. O processo lógico de como depurar um programa não é abordado neste capítulo.

\section{Iniciando o DDD}
O depurador \textbf{ddd} é iniciado com o arquivo executável. O programa deve ser montado e vinculado às opções corretas (conforme observado no capítulo anterior). Por exemplo, usando o programa exemplo anterior, por exemplo, o comando seria:
\begin{verbatim}
ddd primeiro
\end{verbatim}

Ao iniciar o DDD/GDB, algo semelhante à tela, mostrada abaixo, deve ser exibida (com o código-fonte apropriado exibido).

\begin{figure}
\begin{center}
		\includegraphics[width=0.8\linewidth]{imagens/ddd01}
\end{center}
	\caption{Tela inicial do depurador.}
\end{figure}

Se o código não for exibido de maneira semelhante à mostrada acima, as etapas de montagem e vinculação devem ser verificadas. Especificamente, o qualificador \textbf{-g} deve ser incluído nas etapas de montagem e de link.

A ajuda integrada está disponível clicando no item de menu \textbf{Help} (canto superior direito). Os manuais \textbf{DDD} e \textbf{GDB} estão disponíveis nas páginas dos respectivos desenvolvedores. Para sair do DDD/GDB, selecione \textbf{File} $\rightarrow$ \textbf{Exit} (na barra de menu superior).

\subsection{Definições de configuração DDD}
Algumas sugestões adicionais de configurações de \textbf{DDD/GDB} incluem:

\begin{itemize}
	\item \textbf{Edit $\rightarrow$ Preferences $\rightarrow$ General $\rightarrow$ Suppress X Warning}
	\item \textbf{Edit $\rightarrow$ Preferences $\rightarrow$ Source $\rightarrow$ Display Source Line Numbers}
\end{itemize}

Eles não são obrigatórios, mas podem facilitar o uso do depurador. Se definido, as opções serão salvas e lembradas para usos sucessivos do depurador (na mesma máquina).

\section{Execução de programa com DDD}
Para executar o programa, clique no botão \textbf{Run} no menu da ferramenta de comando (mostrado abaixo). Como alternativa, você pode digitar \textbf{run} no prompt \textbf{(gdb)} (janela inferior do console \textbf{GDB}). No entanto, isso executará o programa inteiramente e, quando feito, os resultados serão zerados (e perdidos).

\subsection{Definindo pontos de interrupção}
Para controlar a execução do programa, será necessário definir um \textit{breakpoint} (local de pausa de execução) para pausar o programa em um local selecionado pelo usuário. Isso pode ser feito selecionando o local de origem (linha para parar). Para este exemplo, vamos parar na linha 31.

O ponto de interrupção pode ser feito de três maneiras:\begin{itemize}
	\item Clique com o botão direito no número da linha e selecione: \textit{Set Breakpoint}
	\item No console de comando GDB, no prompt (gdb), digite: \textbf{break last}
	\item No console de comando GDB, no prompt (gdb), digite: \textbf{break 88}
\end{itemize}

No exemplo a seguir, a linha 87 é um rótulo sem instrução. Se um ponto de interrupção for definido no rótulo, ele irá parar na próxima instrução executável (linha 88 neste exemplo).

Quando configurado corretamente, o ícone ``stop'' aparecerá ao lado do número da linha (conforme mostrado no diagrama).

\begin{figure}[h]
	\begin{center}
		\includegraphics[width=0.8\linewidth]{imagens/dddbreakpoint}
	\end{center}
	\caption{Tela do debugger com breakpoint.}
\end{figure}

Os comandos \textbf{DDD/GDB} podem ser digitados na janela inferior (no \textit{prompt (gdb)}) a qualquer momento. Vários pontos de interrupção podem ser definidos, se desejado.

\subsection{Executando Programas}
Uma vez que o depurador é iniciado, a fim de usar efetivamente o depurador, um ponto de interrupção inicial deve ser definido.

Depois que o ponto de interrupção é definido, o comando de execução pode ser executado clicando na janela do menu \textbf{Run} ou digitando \textbf{run} no prompt (gdb). O programa será executado até, \textit{mas não incluindo}, a instrução com a seta verde.

\begin{figure}[h]
	\begin{center}
		\includegraphics[width=0.8\linewidth]{imagens/ddd03}
	\end{center}
	\caption{Tela do debugger com seta verde.}
\end{figure}

O ponto de interrupção é indicado com o sinal de parada à direita (consulte o exemplo anterior). A localização atual é indicada com uma seta verde. Especificamente, a seta verde aponta para a próxima instrução a ser executada. Ou seja, a instrução apontada pela seta verde ainda não foi executada.

\subsubsection{Run/Continue}
Conforme necessário, pontos de interrupção adicionais podem ser definidos. No entanto, clique no comando \textbf{Run} para reiniciar a execução do início e parar no ponto de interrupção inicial.

\begin{figure}[h]
	\begin{center}
		\includegraphics[width=0.12\linewidth]{imagens/ddd04}
	\end{center}
	\caption{Barra de comandos do DDD.}
\end{figure}

Após o comando \textbf{Run} inicial, para continuar para o próximo ponto de interrupção, o comando continue deve ser usado (clicando na janela do menu \textbf{Cont} ou digitando \textbf{cont} no prompt (gdb)). Linhas únicas também podem ser executadas uma linha por vez digitando os comandos passo ou próximo  (clicando em \textbf{Step} ou \textbf{Next} na janela de menu ou digitando \textbf{step} ou \textbf{next} no prompt (gdb)).


\subsubsection{Next/Step}
O comando \textbf{next} executa a próxima instrução. Isso inclui a execução de uma função inteira, se necessário. O comando \textbf{step} executará uma etapa, entrando nas funções se necessário. Para uma única instrução sem função, não há diferença entre os comandos \textbf{next} e \textbf{step}.

\subsection{Exibindo Conteúdo de Registradores}
O método mais simples de ver o conteúdo dos registradores é usar a janela de registradores. A janela de registradores não é exibida por padrão, mas pode ser visualizada selecionando \textbf{Status} $\rightarrow$ \textbf{Registers} (na barra de menu superior). Quando exibida, a janela de registro mostrará o conteúdo do registradores por nome do registrador (coluna da esquerda), em hexadecimal (coluna do meio) e decimal com sinal (coluna da direita).

\begin{figure}[h]
	\begin{center}
		\includegraphics[width=0.5\linewidth]{imagens/ddd05}
	\end{center}
	\caption{Janela de registradores.}
\end{figure}

Dependendo da máquina e da resolução da tela, a janela de registradores pode precisar ser redimensionada para visualizar todo o conteúdo.

A terceira coluna da janela de registradores geralmente mostra a representação decimal de palavras quádruplas, exceto para alguns registradores de propósito especial (\textbf{rbp} e \textbf{rsp}). A representação decimal com quatro palavras com sinais nem sempre é significativa. Por exemplo, se dados não assinados estiverem sendo usados (como endereços), a representação com sinal estará incorreta. Além disso, quando dados de caractere são usados, a representação com sinal não seria significativa.

Por padrão, apenas os registradores inteiros são exibidos. Clicar na caixa ``All registers'' adicionará os registros de ponto flutuante ao \textit{display}. A visualização exigirá rolar para baixo na janela de registradores.

\subsection{Resumo de comandos DDD/GDB}
A tabela a seguir fornece um pequeno subconjunto dos comandos \textbf{DDD} mais comuns. Quando digitados, a maioria dos comandos pode ser abreviada. Por exemplo, \textbf{quit} pode ser abreviado como \textbf{q}. O comando e a abreviatura são mostrados na Tabela \ref{tab:comandosddd}.

\begin{table}[ht]
	\begin{center}
	\resizebox{0.6\textwidth}{!}{%
	\begin{tabular}{|ll|}
	\hline
	\rowcolor[HTML]{C0C0C0}
	\textbf{Comando} & \textbf{Descrição} \\ \hline
	\textbf{quit |q} & Sai do depurador. \\ \hline
	\textbf{break <label/addr>}  & Define um \textit{breakpoint} (ponto de parada) \\ 
	\textbf{| b <label/addr>} & em \textbf{<label>} ou	\textbf{<address>}. \\\hline	
	\textbf{run <args> | r <args>} & Execute o programa (até o primeiro \\
	& ponto de parada).\\ \hline
	\textbf{continue | c} & Continue a execução (até o próximo\\
	&  ponto de parada).\\ \hline
	\textbf{continue <n> | c <n>} & Continue a execução (até o próximo\\
	& ponto de parada), ignorando\\
	& n-1 passagens do ponto de parada.\\
	& Isso pode ser usado para chegar rapida-\\
	&mente à n-ésima iteração de um loop.\\ \hline
	\textbf{step | s} & Passe para a próxima instrução\\
	& (ou seja, siga as chamadas\\
	& de função/procedimento)\\ \hline
	\textbf{next | n }& Próxima instrução (não entra em\\
	& funções/chamadas de procedimento).\\ \hline
	\textbf{F3} & Reinicie o programa (e pare no\\
	& primeiro ponto de parada).\\\hline
	\textbf{where} & Ativação atual (profundidade\\
	& da chamada).\\ \hline
	\textbf{x/<n><f><u> \$rsp} & Examine o conteúdo da pilha.\\ \hline
	\textbf{x\/<n><f><u> \&<var>} & Examine a localização da memória\\
	& <var>. <n> número de locais\\
	& a serem exibidos, 1 é o padrão.\\
	&\begin{tabular}{ll}
		\textbf{<f> formato}: &d - decimal (com sinal)\\
	    &x - hex\\
	    &u - decimal (sem sinal)\\
	    &c - caractere\\
	    &s - string\\
	    &f - ponto flutuante\\	
	\end{tabular}\\
	& \textbf{<u> tamanho da unidade}: \\
	&\begin{tabular}{ll}
		&b - byte (8 bits)\\
	\hspace{2.55cm}	& h - halfword (16-bits)\\
		& w - word (32-bits)\\
        & g - gigante (64-bits) \\
	\end{tabular}\\	\hline
	source <filename> & Leia os comandos do arquivo \textbf{<filename>}.\\\hline
	\textbf{set logging file <filename>}& Defina o arquivo de registro como\\
	& \textbf{<filename>}, o padrão é gdb.txt.\\\hline
	\textbf{set logging on} & Ative o registro (para um arquivo). \\\hline
	\textbf{set logging off}	& Desative o registro (para um arquivo). \\\hline
	\textbf{set logging overwrite}&Quando o registro (para um\\
	& arquivo) estiver ativado, substitua o\\
	&arquivo de registro anterior (se houver). \\\hline
\end{tabular}}
\end{center}
	\caption{Comandos DDD.}
	\label{tab:comandosddd}
\end{table}

Mais informações podem ser obtidas por meio do recurso de ajuda integrado ou da documentação no site do \textbf{ddd} (referenciado no Capítulo \ref{cap1}.

\subsubsection{Comandos DDD/GDB, exemplos}
Por exemplo, dadas as declarações de dados abaixo:

\lstinputlisting{source/exemplo1.asm}

Assumindo dados com sinal, os comandos para examinar as instruções de memória seriam os seguintes:
\begin{verbatim}
x/db &bnum1
x/dh &wnum2
x/dw &dnum3
x/dg &qnum
x/s &class
x/f &twopi
\end{verbatim}

Se um comando de exame de memória (\textit{memory dump}) inadequado for usado (ou seja, tamanho incorreto), não haverá mensagem de erro e o depurador exibirá o que foi solicitado (mesmo que não faça sentido). O exame das variáveis exigirá o uso do comando de exame de memória  apropriado com base nas declarações de dados. As opções adicionais podem ser acessadas no menu na parte superior da tela.

Para exibir um \textit{array} em \textbf{DDD}, o comando básico de exame de memória é usado.
\begin{center}
	\textbf{x/<n><f><u> \&<variable>}
\end{center}

Por exemplo, assumindo a declaração de:
\begin{lstlisting}
list1	dd	100001, -100002, 100003, 100004, 100005
\end{lstlisting}

O comando de exame de memória seria o seguinte:\\
\textbf{x/5dw \&list1}

Onde o \textbf{5} é o comprimento do array. O \textbf{d} indica dados com sinal (\textbf{u} teria sido dados sem sinal). O \textbf{w} indica dados de tamanho de \textbf{32 bits} (que é o que a definição \textbf{dd}, \textit{define double}, declara no arquivo de origem). O \textbf{\&list1} se refere ao endereço da variável. Observe que o endereço aponta para o primeiro elemento (e apenas o primeiro elemento). Como tal, é possível exibir menos ou mais elementos que são realmente declarados no array.

O comando básico para examinar a memória pode ser usado diretamente com um endereço de memória (em oposição a um nome de variável). Por exemplo:\\
\textbf{x/dw 0x600d44}

Os endereços são normalmente exibidos em hexadecimal, portanto, um \textbf{0x} seria necessário para inserir o endereço hexadecimal diretamente, conforme mostrado.

\subsection{Exibindo o conteúdo da pilha}
Existem algumas ocasiões em que a exibição do conteúdo da pilha pode ser útil. A pilha é normalmente composta de elementos não sinalizados de 64 bits. O comando para exame de memória é usado, porém o endereço está no registrador \textbf{rsp} (não é um nome de variável). O comando de exame de memória para exibir o topo atual da pilha seria o seguinte:\\
\textbf{x/ug \$rsp}

O comando de exame de memória para exibir os 6 itens do topo na pilha seria o seguinte:\\
\textbf{x/6ug \$rsp}

Devido à implementação da pilha, o primeiro item mostrado sempre será o topo atual da pilha.

\subsection{Arquivo de comandos do depurador}
Como os comandos de exibição de dados devem estar corretos (já que não há erro), pode ser entediante. Para ajudar a reduzir os erros, a execução correta e o comando de exibição podem ser armazenados em um arquivo de texto. O depurador pode então ler os comandos do arquivo (em vez de digitá-los manualmente). Embora os resultados sejam normalmente exibidos na tela, os resultados podem ser redirecionados para um arquivo de saída. Isso pode ser útil para uma revisão fácil.

Por exemplo, alguns comandos típicos do depurador para definir o ponto de interrupção, executar o programa, exibir algumas variáveis e redirecionar a saída para um arquivo de log podem ser os seguintes:

\lstinputlisting{source/script1.txt}
\textit{Nota 1}; este exemplo assume que um rótulo '\textbf{last}' é definido no programa de origem (como é feito no programa de exemplo).

\textit{Nota 2}; este exemplo sai do depurador. Se isso não for desejado, o comando 'quit' pode ser removido. Ao sair do arquivo de entrada, o depurador pode solicitar a confirmação do usuário da saída (sim ou não).

Esses comandos devem ser colocados em um arquivo (como \textit{gdbIn.tx}t), para que possam ser lidos a partir do depurador.

\subsubsection{Arquivo de comandos do depurador (interativo)}
O comando do depurador para ler um arquivo é ''source <filename>''. Por exemplo, se o nome do arquivo de comando for gdbIn.txt,\\
	\textbf{(gdb) source gdbIn.txt}

Com base nos comandos acima, a saída será colocada no arquivo \textit{out.txt}. O nome do arquivo de saída pode ser alterado conforme desejado.

Cada programa exigirá um conjunto de comandos de entrada personalizado com base nas variáveis e tamanhos associados. O arquivo de comandos de entrada do depurador só será útil quando o programa estiver quase funcionando. Travamentos do programa e outros erros mais significativos exigirão depuração interativa para determinar o erro específico.

\subsubsection{Arquivo de comandos do depurador (não interativo)}
É possível obter o arquivo de saída diretamente sem uma sessão \textbf{DDD} interativa. O comando a seguir, inserido na linha de comando, executará o comando no arquivo de entrada do programa fornecido, criará o arquivo de saída e sairá do programa.\\
\textbf{gbd <gdbIn.txt prog}\\
Que criará o arquivo de saída (conforme especificado no arquivo \textit{gdbIn.txt}) e sairá do depurador. Esta é a opção mais rápida para obter o arquivo de saída final para um programa em funcionamento. Novamente, isso só seria útil se o programa estiver funcionando ou muito perto de funcionar corretamente.

\section{Exercícios}
Abaixo estão algumas perguntas baseadas neste capítulo.

\subsection{Questionário}
\begin{enumerate}
	\item Como o depurador é iniciado (na linha de comando)?
	\item Qual opção é necessária durante a etapa de montagem e vinculação para garantir o	programa ser facilmente depurado.
	\item O que o comando \textbf{run} faz especificamente?
	\item O que o comando \textbf{continue} faz especificamente?
	\item Como a janela de registradores é exibida?
	\item Existem três colunas na janela de registradores. O primeiro mostra o registrador.	O que as outras duas colunas mostram?
	\item Uma vez que o depurador é iniciado, como o usuário pode sair?
	\item Descreva como um ponto de parada é definido (várias maneiras).
	\item Qual é o comando debugger para ler os comandos debugger de um arquivo?
	\item Quando o DDD mostra uma seta verde apontando para uma instrução, o que isso significa?
	\item Forneça o comando debugger para exibir cada uma das seguintes variáveis em decimal.
	\begin{enumerate}
		\item \textbf{bVar1} (variável de tamanho de byte)
		\item \textbf{wVar1} (variável de tamanho de palavra)
		\item \textbf{dVar1} (variável de tamanho de palavra dupla)
		\item \textbf{qVar1} (variável de tamanho de palavra quádrupla)
		\item \textbf{bArr1} (matriz de 30 elementos de bytes)
		\item \textbf{wArr1} (matriz de palavras com 50 elementos)
		\item \textbf{dArr1} (matriz de 75 elementos de palavras duplas)
	\end{enumerate}
    \item Forneça o comando debugger para exibir cada uma das seguintes variáveis em formato hexadecimal.
    \begin{enumerate}
    	\item \textbf{bVar1} (variável de tamanho de byte)
    	\item \textbf{wVar1} (variável de tamanho de palavra)
    	\item \textbf{dVar1} (variável de tamanho de palavra dupla)
    	\item \textbf{qVar1} (variável de tamanho de palavra quádrupla)
    	\item \textbf{bArr1} (matriz de 30 elementos de bytes)
    	\item \textbf{wArr1} (matriz de palavras com 50 elementos)
    	\item \textbf{dArr1} (matriz de 75 elementos de palavras duplas)
    \end{enumerate}
    \item  Qual é o comando do depurador para exibir o valor no topo da pilha atual?
    \item Qual é o comando debugger para exibir cinco (5) valores no topo da pilha atual?
\end{enumerate}

\subsection{Projetos sugeridos}
Abaixo estão alguns projetos sugeridos com base neste capítulo.
\begin{enumerate}
	\item Digite o programa de exemplo do Capítulo \ref{cap4:formato}, Formato do programa. Monte e conecte o programa conforme descrito no Capítulo \ref{Ferramentasdedesenvolvimento}, Ferramentas de desenvolvimento. Execute o depurador conforme observado neste capítulo. Defina um ponto de interrupção no último rótulo e execute o programa (para esse ponto de interrupção). Verifique interativamente se os cálculos realizados resultaram nos valores corretos. Isso exigirá a digitação dos comandos de memória de exame do depurador apropriados (com base no tamanho da variável).
	\item Depois de concluir o problema anterior, crie um arquivo de entrada do depurador que definirá o envio da saída para um arquivo de texto, definirá um ponto de interrupção, executará o programa e exibirá os resultados para cada variável (com base no tamanho de variável apropriado). Execute o depurador e leia o arquivo de origem. Verifique se o arquivo de entrada funcionou corretamente e se os cálculos do programa estão corretos com base nos resultados mostrados no arquivo de saída.
	\item Crie um arquivo de script de montagem e link, conforme descrito no Capítulo \ref{Ferramentasdedesenvolvimento}, Ferramentas de desenvolvimento. Use o script para montar e linkar o programa. Certifique-se de que o script monte e link corretamente.
\end{enumerate}



\chapter{Visão Geral do Conjunto de Instruções}

\begin{chapquote}{Autor desconhecido}
``Por que os livros de matemática são tristes?
Porque eles têm muitos problemas.''
\end{chapquote}

Este capítulo fornece uma visão geral básica de um subconjunto simples do conjunto de instruções x86-64 com foco nas operações inteiras. Isso abrangerá apenas o subconjunto de instruções necessárias para os tópicos e programas discutidos no escopo deste texto. Isso excluirá algumas das instruções mais avançadas e instruções de modo restrito. Para obter uma lista completa de todas as instruções do processador, consulte as referências listadas no Capítulo \ref{cap1}.

As instruções são apresentadas na seguinte ordem:
\begin{itemize}
	\item Movimentação de dados
	\item Instruções de conversão
	\item Instruções aritméticas
	\item Instruções lógicas
	\item Instruções de controle
\end{itemize}

As instruções para chamadas de funções são discutidas no capítulo no capítulo \ref{cap12}, Funções.

Uma lista completa das instruções abordadas neste texto está localizada no Apêndice \ref{apendiceB} para referência.

\section{Convenções notacionais}
Esta seção resume a notação usada neste texto, que é bastante comum na literatura técnica. Em geral, uma instrução consiste na própria instrução ou operação (isto é, add, sub, mul, etc.) e nos \textbf{operandos}. Os operandos se referem de onde os dados (a serem operados) são provenientes e/ou onde o resultado deve ser colocado.

\subsection{Notação de Operando}
A tabela a seguir resume as convenções de notação usadas no restante do documento.

{\centering
	\begin{table}[ht]
		\centering
		\begin{tabular}{|p{3cm}|p{9cm}|}
			\hline
			\rowcolor[HTML]{C0C0C0} 
			{\color[HTML]{000000} } Notação de Operando & {\color[HTML]{000000} }Descrição \\ \hline
			\textbf{<reg>}& Operando registrador. O operando deve ser um registrador. \\ \hline
			\textbf{<reg8>,<reg16>, <reg32>,<reg64>} & Registrador operando com requisito de tamanho específico. Por exemplo, \textbf{reg8} significa apenas um registro de tamanho de byte (por exemplo, \textbf{al}, \textbf{bl}, etc.) e \textbf{reg32} significa um registro de tamanho de palavra dupla (por exemplo, \textbf{eax}, \textbf{ebx}, etc.).\\ \hline
			\textbf{<dest>} & Operando de destino. O operando pode ser um registro ou memória. Por ser um operando de destino, o conteúdo será sobrescrito pelo novo resultado (com base nas instruções específicas).\\ \hline
			\textbf{<RXdest>} & Operando de registro de destino de ponto flutuante. O operando deve ser um registrador de ponto flutuante. Por ser um operando de destino, o conteúdo será sobrescrito pelo novo resultado (com base nas instruções específicas).\\ \hline
			\textbf{<src>} & Operando de origem. O valor do operando é inalterado após a instrução.\\ \hline
			\textbf{<imm>} & Valor imediato. Pode ser especificado em decimal, hexadecimal, octal ou binário.\\ \hline
			\textbf{<mem>} & Localização da memória. Pode ser um nome de variável ou uma referência indireta (ou seja, um endereço de memória).\\ \hline
			\textbf{<op> ou <operand>} & Operando, registrador ou memória.\\ \hline
			\textbf{<op8>,<op16>, <op32>,<op64>} & Operando, registrador ou memória, com requisitos de tamanho específico. Por exemplo, \textbf{op8} significa apenas um operando de tamanho de byte e \textbf{reg32} significa apenas um operando de tamanho de palavra dupla.\\ \hline
			\textbf{<label>} & Rótulo do programa.\\ \hline
		\end{tabular}
		\caption{}
		\label{notacao}
	\end{table}
}


\chapter{Modos de Endereçamento}
\label{cap:ModosDeEndereçamento}

\chapter{Pilha de Processos}
\label{cap:pilhadeprocessos}

\chapter{Desenvolvimento de Programas}
\label{cap:DesenvolvimentoDeProgramas}

\chapter{Macros}
\label{cap:Macros}

\chapter{Funções}
\label{cap:Funções}


\end{document}
